\documentclass[11pt]{article}

\usepackage{amsmath,amssymb,amsthm}
\usepackage{mathtools}
\usepackage{geometry}
\usepackage{hyperref}
\usepackage{natbib}
\usepackage{graphicx}
\usepackage{booktabs}
\usepackage{microtype}
\usepackage{enumitem}
\usepackage{xcolor}
\usepackage{float}

\geometry{margin=1in}

\newtheorem{theorem}{Theorem}[section]
\newtheorem{proposition}[theorem]{Proposition}
\newtheorem{lemma}[theorem]{Lemma}
\newtheorem{corollary}[theorem]{Corollary}
\newtheorem{conjecture}[theorem]{Conjecture}
\theoremstyle{definition}
\newtheorem{definition}[theorem]{Definition}
\newtheorem{example}[theorem]{Example}
\theoremstyle{remark}
\newtheorem{remark}[theorem]{Remark}

\newcommand{\E}{\mathbb{E}}
\newcommand{\PP}{\mathbb{P}}
\newcommand{\R}{\mathbb{R}}
\newcommand{\Z}{\mathbb{Z}}
\newcommand{\N}{\mathbb{N}}
\newcommand{\ind}{\mathbf{1}}
\newcommand{\dd}{\,\mathrm{d}}
\newcommand{\Fees}{\mathrm{Fees}}
\newcommand{\IL}{\mathrm{IL}}

\title{Finite-Tick Corrections for Concentrated Liquidity:\\
Non-Universality, Misalignment, and Optimal Range Design}
\author{}
\date{\today}

\begin{document}
\maketitle

\begin{abstract}
Echenim, Gobet, and Maurice (2025) derived asymptotic formulas for
liquidity provider (LP) fee revenue in Uniswap~v3, proving convergence
to an integral of the price local time as the tick spacing
$\gamma \to 0$. In practice, $\gamma$ is fixed by the pool's fee tier
and can be as large as $2\%$. We derive three sets of results on the
finite-tick regime. \textbf{(1) Non-universality of $\kappa_1$.} Under a
general constant-drift diffusion $dX_t = \mu\dd t + \sigma\dd W_t$, the
first-order correction coefficient is $\kappa_1 = 1/2 + \mu/(2\sigma^2)$,
which equals $1/4$ \emph{only} for risk-neutral geometric Brownian
motion ($\mu = -\sigma^2/2$). Under the physical measure with expected
return $r > 0$, the correction is $\kappa_1^P = 1/4 + r/(2\sigma^2)$,
which for typical crypto parameters is $2$--$3$ times the risk-neutral
value. The result decomposes into a prefactor contribution and an
Euler--Maclaurin boundary term, and the cancellation $f'(0) = 0$ that
produces ``universality'' is a precise algebraic artifact of the
risk-neutral GBM drift. \textbf{(2) Misaligned-range correction.} When
the LP range boundary falls between two tick levels with fractional
offset $\delta$, the boundary tick earns a fraction
$R(\delta) = 2(1-\delta)/(2-\delta)$ of the interior tick's fee. Unlike
the classical Broadie--Glasserman--Kou correction for temporally
discretized barriers, this spatial correction is a rational function of
$\delta$ and does not involve the Wiener--Hopf constant
$\zeta(1/2)/\sqrt{2\pi}$, because the continuous price process crosses
every tick level with probability one. \textbf{(3) Optimal range width.}
The finite-tick correction shifts the fee-maximizing range width downward
by $1$--$2$ ticks relative to the EGM asymptotic, with the largest
effect in the $0.05\%$ pool. We validate all results against
semi-analytical computations.
\end{abstract}

%% ===================================================================
\section{Introduction}\label{sec:intro}

%% -------------------------------------------------------------------
\subsection{Concentrated liquidity in automated market makers}

Automated Market Makers (AMMs) enable decentralized exchange of
crypto-assets without a traditional order book. The constant-product
formula $x \cdot y = k$, introduced by Uniswap~v1~\citep{Adams2021},
requires liquidity providers (LPs) to spread capital across the entire
price range $[0, \infty)$, resulting in low capital efficiency.
Uniswap~v3~\citep{Adams2021} addressed this through
\emph{concentrated liquidity}: LPs restrict their capital to a chosen
price range $[p_\ell, p_u]$, earning swap fees only when the market
price falls within this range. The concentration amplifies fee income
per unit of capital but exposes the LP to \emph{impermanent loss} (IL)
when the price exits the range.

The tick mechanism discretizes the price space into levels spaced
multiplicatively by $(1 + \gamma)$, where $\gamma$ depends on the
pool's fee tier:
\begin{table}[h]
\centering
\caption{Uniswap v3 fee tiers and tick spacings.}
\label{tab:fee_tiers}
\begin{tabular}{cccc}
\toprule
Fee tier $\phi$ & Tick spacing $\delta_\pi$ & $\gamma$ (approx.)
& Typical use case \\
\midrule
$0.01\%$ & $1$ & $0.01\%$ & Stablecoins \\
$0.05\%$ & $10$ & $0.10\%$ & Correlated pairs \\
$0.30\%$ & $60$ & $0.60\%$ & Most pairs \\
$1.00\%$ & $200$ & $2.02\%$ & Exotic/volatile \\
\bottomrule
\end{tabular}
\end{table}

\noindent
The tick spacing is coarser for higher fee tiers, reflecting the
intuition that volatile pairs require wider spreads. Uniswap~v4
retains the same tick mechanism with identical fee tier structure, so
all results in this paper apply equally to both versions.

The discreteness of the tick grid introduces a systematic gap between
the continuous-limit fee formulas used for LP strategy optimization and
the actual fee revenue earned on-chain. Quantifying this gap is the
central problem of this paper.

%% -------------------------------------------------------------------
\subsection{The EGM asymptotic framework}

\citet{EGM2025} established a rigorous asymptotic framework for LP fee
revenue. Their main result (Theorem~4.1) shows that as
$\gamma \to 0$, the renormalized fee revenue
$\gamma \cdot \Fees$ converges in probability to an integral
involving the local time of the log-price process:
\begin{equation}\label{eq:EGM_limit}
\gamma \cdot \Fees \xrightarrow{\PP} \frac{\phi}{1 - \phi}
\int_{p_\ell}^{p_u} A^b_T(p)\, b^{-3/2}\dd b,
\end{equation}
where $A^b_T(p)$ is the occupation functional of the price process at
level~$b$, which under continuous diffusion models reduces to the
semimartingale local time $L^{\log b}_T(X)/b$ of the log-price. Under
geometric Brownian motion (GBM), they obtained an explicit
expected-fee formula (Theorem~4.5). These results hold in the limit
$\gamma \to 0$; the finite-$\gamma$ regime was left unexplored.

The practical relevance of finite-$\gamma$ corrections is immediate:
for the $1\%$ fee tier, $\gamma \approx 0.02$, and for a 20-tick LP
range, the correction to the EGM asymptotic can exceed $1.5\%$ of
total fee revenue. For an LP deploying \$10M of capital, this
corresponds to approximately \$15{,}000 per year in systematic
estimation error.

%% -------------------------------------------------------------------
\subsection{Our contributions}

We derive three sets of results that substantially extend the EGM
framework to the finite-tick regime.

\paragraph{Contribution 1: Non-universality of $\kappa_1$.}
Under any constant-drift diffusion
$dX_t = \mu\dd t + \sigma\dd W_t$, we prove that the first-order
correction coefficient is
\begin{equation}\label{eq:kappa1_intro}
\kappa_1(\mu, \sigma) = \frac{1}{2} + \frac{\mu}{2\sigma^2},
\end{equation}
so that $\Psi(\gamma)/F_0^Y(\gamma) = 1 + \kappa_1 \gamma +
O(\gamma^2)$. The value $\kappa_1 = 1/4$ holds \emph{only} for the
risk-neutral GBM drift $\mu = -\sigma^2/2$. Under the physical measure
with expected excess return $r > 0$, the correction becomes
$\kappa_1^P = 1/4 + r/(2\sigma^2)$, which for Bitcoin
($r \approx 0.5$, $\sigma \approx 0.8$) gives
$\kappa_1^P \approx 0.64$---more than $2.5$ times the risk-neutral
value. The result decomposes into a prefactor term (from the nonlinear
expansion of the fee-per-crossing formula) and an Euler--Maclaurin
boundary term (from the derivative $f'(0)$ of the weight function at
the origin). The precise cancellation $f'(0) = 0$ that produces
``universality'' is shown to be an algebraic artifact of the
risk-neutral GBM drift.

\paragraph{Contribution 2: Misaligned-range correction.}
When the LP range boundary does not align with the tick grid, the
boundary tick earns a reduced fraction of the interior tick's fee. We
derive the explicit formula
\[
R(\delta) = \frac{2(1-\delta)}{2-\delta}
\]
for the boundary tick efficiency ratio, where $\delta \in [0,1)$ is
the fractional offset. This correction is a smooth rational function
that interpolates from $R(0) = 1$ (aligned) to $R(1^-) = 0$
(boundary at the tick level). We prove that, unlike the classical
Broadie--Glasserman--Kou (BGK) correction for temporally discretized
barriers, the spatial misalignment correction does \emph{not} involve
the Wiener--Hopf constant $\beta = \zeta(1/2)/\sqrt{2\pi}$. The
fundamental reason is that a continuous price process crosses every
spatial level with probability one (no overshoot), whereas discrete
temporal monitoring allows the process to jump past a barrier between
observation dates.

\paragraph{Contribution 3: Optimal range width.}
We derive the fee-maximizing range width $N^*(\gamma)$ under the
corrected fee formula and show that the second-order correction
$-N^2\gamma^2/6$ acts as an additional penalty on wide ranges,
shifting the optimum downward by $1$--$2$ ticks relative to the EGM
asymptotic. The correction is most significant for the $0.05\%$ pool
($\gamma = 0.001$) where optimal ranges are widest. For a \$1M LP
position in this pool, the corrected range design improves daily fee
revenue by approximately \$200.

%% -------------------------------------------------------------------
\subsection{Mathematical structure: two sources of correction}

The finite-tick correction has a clean two-source decomposition that
illuminates its mathematical structure:
\begin{enumerate}[label=(\roman*)]
\item \textbf{Prefactor correction} $\kappa_1^{\psi}$: The fee per
crossing involves the nonlinear function
$c(h) = (\gamma/h)(e^{h/2} - 1) \cdot S(h)/S(2h)$, where $S$ is the
scale function. Taylor expanding $c(h)$ around $h = 0$ gives a
first-order contribution $\kappa_1^{\psi} = 3/4 + \mu/\sigma^2$.
\item \textbf{Euler--Maclaurin correction} $\kappa_1^{EM}$: The
discrete sum $\sum h f(kh)$ over the tick grid approximates the
integral $\int f(x)\dd x$ with an error governed by the
Euler--Maclaurin formula. When $f'(0) \ne 0$, the leading boundary
term contributes
$\kappa_1^{EM} = -1/4 - \mu/(2\sigma^2)$.
\end{enumerate}
The total correction is
$\kappa_1 = \kappa_1^{\psi} + \kappa_1^{EM} = 1/2 + \mu/(2\sigma^2)$.
Under risk-neutral GBM ($\mu = -\sigma^2/2$), the Euler--Maclaurin
contribution simplifies to $\kappa_1^{EM} = 0$ (because $f'(0) = 0$),
leaving only $\kappa_1^{\psi} = 1/4$. This cancellation is
\emph{specific} to the risk-neutral drift; any other drift breaks it.

%% -------------------------------------------------------------------
\subsection{Connection to discrete barrier corrections}

The classical result of \citet{BGK1997,BGK1999} shows that a
discretely monitored barrier option price can be corrected by shifting
the barrier by $\beta\sigma\sqrt{\Delta t}$, where
$\beta = \zeta(1/2)/\sqrt{2\pi} \approx -0.5826$ arises from the
Wiener--Hopf factorization of the random walk maximum's distribution.
\citet{Horfelt2003} extended this to second order, and
\citet{HowisonSteinberg2007} provided an alternative derivation via
matched asymptotic expansions. \citet{FengLinetsky2008} further
extended the framework to exponential L\'evy processes using
Hilbert transforms.

In our setting, the tick grid forms a \emph{spatial} discretization,
not a temporal one. The price process is continuous in time and crosses
each tick level with probability one---there is no random walk
overshoot, so the Spitzer--Baxter machinery underlying BGK does not
apply. Our corrections are instead algebraic (from the fee formula
expansion) and analytic (from the Euler--Maclaurin formula applied to
the crossing-count sum). For misaligned range boundaries, the
correction is a rational function of the fractional offset, not an
irrational constant involving $\zeta(1/2)$.

This distinction between spatial and temporal discretization
corrections appears to be new in the AMM literature, though it is
implicit in the matched asymptotics framework of
\citet{HowisonSteinberg2007}.

%% -------------------------------------------------------------------
\subsection{Related work}

The qualitative connection between AMM positions and barrier options
was noted by \citet{Paone2025}. The technical backbone of the EGM
proof builds on the local time approximation results of
\citet{GobetLandon2014} and \citet{GobetStazhynski2020}.
\citet{ClarkHassani2023} analyzed LP returns under GBM using option
pricing techniques. \citet{FukasawaMajoWorhacz2023} studied optimal
concentrated liquidity provision as a continuous-time stochastic
control problem. \citet{CartaFabrizio2024} derived fee revenue
formulas for general semimartingale price processes using
multiplicative Doob decompositions. To our knowledge, no prior work
derives a finite-tick correction for AMM fee revenue, nor identifies
the non-universality of the first-order correction coefficient.

%% -------------------------------------------------------------------
\subsection{Organization}

Section~\ref{sec:setup} introduces the model setup and notation.
Section~\ref{sec:general_correction} states and proves the general
first-order correction formula for arbitrary constant drift.
Section~\ref{sec:cancellation} provides a structural analysis of the
$f'(0) = 0$ cancellation and its consequences.
Section~\ref{sec:gbm} specializes to GBM and derives the full
two-term correction including the second-order range-dependent term.
Section~\ref{sec:misaligned} treats misaligned range boundaries.
Section~\ref{sec:optimal} derives the optimal range width under the
corrected formula. Section~\ref{sec:numerics} presents numerical
validation. Section~\ref{sec:discussion} discusses economic
implications and extensions. Section~\ref{sec:conclusion} concludes.
Detailed calculations are collected in the Appendix.

%% ===================================================================
\section{Setup}\label{sec:setup}

%% -------------------------------------------------------------------
\subsection{Price dynamics and the log-price process}

Let $p_t$ denote the exchange rate (price of asset~$X$ in units of
asset~$Y$) at time $t \in [0, T]$. Under a general constant-drift
diffusion model, the log-price $X_t = \log p_t$ satisfies
\begin{equation}\label{eq:general_sde}
dX_t = \mu\dd t + \sigma\dd W_t, \qquad X_0 = x_0,
\end{equation}
with constant drift $\mu \in \R$ and volatility $\sigma > 0$. This
encompasses several important cases:
\begin{itemize}[nosep]
\item \textbf{Risk-neutral GBM}: $\mu = -\sigma^2/2$, corresponding
to $dp_t/p_t = \sigma\dd W_t$ under $\PP^*$.
\item \textbf{Physical-measure GBM}: $\mu = r - \sigma^2/2$, where
$r$ is the expected (log-)return.
\item \textbf{Arithmetic Brownian motion (ABM)}: $\mu = 0$,
$X_t = x_0 + \sigma W_t$, modeling a driftless log-price.
\item \textbf{General drift}: $\mu$ arbitrary, e.g.,
$\mu > 0$ for trending assets, $\mu < 0$ for deflationary tokens.
\end{itemize}

%% -------------------------------------------------------------------
\subsection{Tick grid and LP positions}

The tick spacing parameter $\gamma > 0$ induces a log-price grid with
spacing
\begin{equation}\label{eq:h_def}
h = \log(1 + \gamma) = \gamma - \frac{\gamma^2}{2} +
\frac{\gamma^3}{3} - \cdots.
\end{equation}
The grid levels are $\{kh : k \in \Z\}$. An LP deposits unit
liquidity on $N$ consecutive ticks centered at $x_0 = 0$ (without
loss of generality), giving a range
\begin{equation}\label{eq:range}
R_\gamma = \Bigl[-\frac{Nh}{2},\, \frac{Nh}{2}\Bigr]
\end{equation}
in log-price space. We assume $N$ is even throughout.

%% -------------------------------------------------------------------
\subsection{Scale function and crossing probabilities}\label{sec:scale}

The scale function of the diffusion \eqref{eq:general_sde} is
\begin{equation}\label{eq:scale_fn}
S(x) = \int_0^x \exp\Bigl(-\frac{2\mu}{\sigma^2}\,y\Bigr)\dd y
= \frac{1}{\alpha}\bigl(e^{\alpha x} - 1\bigr),
\end{equation}
where we define the \emph{drift parameter}
\begin{equation}\label{eq:alpha_def}
\alpha \coloneqq -\frac{2\mu}{\sigma^2}.
\end{equation}
For risk-neutral GBM, $\alpha = 1$; for ABM, $\alpha = 0$ (and
$S(x) = x$ by continuity). The gambler's ruin formula gives the
probability of an upcrossing from a grid level:

\begin{lemma}[Crossing probability]\label{lem:p_up}
Starting at grid level $kh$, the probability of reaching
$(k{+}1)h$ before $(k{-}1)h$ under the diffusion
\eqref{eq:general_sde} is
\begin{equation}\label{eq:p_up}
p_\uparrow(h) = \frac{S(h)}{S(2h)}
= \frac{e^{\alpha h} - 1}{e^{2\alpha h} - 1}
= \frac{1}{1 + e^{\alpha h}}.
\end{equation}
In particular:
\begin{align}
p_\uparrow(h) &= \frac{1}{2} - \frac{\alpha h}{4} + O(h^2),
\label{eq:p_up_expand}
\end{align}
so the upcrossing probability is less (resp.\ greater) than $1/2$
when $\alpha > 0$ (resp.\ $\alpha < 0$), i.e., when the drift is
negative (resp.\ positive).
\end{lemma}
\begin{proof}
By the standard exit problem for Brownian motion with drift, the
probability of reaching level $a$ before level $-b$ starting from
$0$ is $S(b)/S(a+b)$ when $S$ is the scale function. Setting
$a = b = h$ and using the explicit form
\eqref{eq:scale_fn} gives
\[
p_\uparrow = \frac{S(h)}{S(2h)}
= \frac{e^{\alpha h} - 1}{e^{2\alpha h} - 1}
= \frac{1}{1 + e^{\alpha h}}.
\]
The last equality follows from the factorization
$e^{2\alpha h} - 1 = (e^{\alpha h} - 1)(e^{\alpha h} + 1)$.
For the Taylor expansion, write $e^{\alpha h} = 1 + \alpha h +
(\alpha h)^2/2 + O(h^3)$, so
$1 + e^{\alpha h} = 2 + \alpha h + O(h^2)$, and
$1/(2 + \alpha h + O(h^2)) = 1/2 - \alpha h/4 + O(h^2)$.
\end{proof}

%% -------------------------------------------------------------------
\subsection{Local time and occupation density}

The semimartingale local time $L^x_T = L^x_T(X)$ of the process $X$
at level $x$ up to time $T$ is defined via the occupation times
formula:
\begin{equation}\label{eq:occupation}
\int_0^T g(X_s)\dd\langle X\rangle_s
= \int_{-\infty}^{\infty} g(x)\,L^x_T\dd x
\end{equation}
for all measurable $g$. Since $\langle X\rangle_t = \sigma^2 t$, this
is equivalent to
$\int_0^T g(X_s)\dd s
= \sigma^{-2}\int g(x)\,L^x_T\dd x$.

The expected local time has the explicit representation
\begin{equation}\label{eq:expected_LT}
\E[L^x_T] = \sigma^2 \int_0^T p_X(t, 0, x)\dd t,
\end{equation}
where $p_X(t, 0, x)$ is the transition density of $X$:
\begin{equation}\label{eq:transition}
p_X(t, 0, x) = \frac{1}{\sigma\sqrt{2\pi t}}
\exp\Bigl(-\frac{(x - \mu t)^2}{2\sigma^2 t}\Bigr).
\end{equation}

%% -------------------------------------------------------------------
\subsection{Fee structure and the crossing-count representation}

Following \citet{EGM2025}, the $Y$-token fee accumulated over
$[0, T]$ from upcrossings at tick level $kh$ is
\begin{equation}\label{eq:fee_per_cross}
\text{fee per upcrossing at } kh
= \pi_k \cdot \bigl((1{+}\gamma)^{1/2} - 1\bigr)
\cdot \frac{\phi}{1 - \phi},
\end{equation}
where $\pi_k = e^{kh/2}$ is the square-root price at tick level $kh$.
The factor $(1+\gamma)^{1/2} - 1$ is the relative change in the
virtual reserve $\sqrt{p}$ across one tick.

\begin{definition}[Prefactor function]\label{def:prefactor}
Define the crossing-count coefficient
\begin{equation}\label{eq:c_def}
c(h) \coloneqq \frac{\gamma}{h} \cdot
\frac{e^{h/2} - 1}{1 + e^{\alpha h}},
\end{equation}
and the normalized prefactor
\begin{equation}\label{eq:psi_def}
\psi(h) \coloneqq \frac{4c(h)}{h}.
\end{equation}
\end{definition}

The expected renormalized $Y$-fee is then
\begin{equation}\label{eq:Psi_def}
\Psi(\gamma) = \frac{\phi}{1 - \phi}
\sum_{k=-N/2}^{N/2-1} c(h) \cdot e^{kh/2} \cdot \E[L^{kh}_T].
\end{equation}

%% -------------------------------------------------------------------
\subsection{The EGM asymptotic formula}

The $Y$-fee component of the EGM limit for the range $R_\gamma$ is
\begin{equation}\label{eq:F0Y}
F_0^Y(\gamma) = \frac{\phi}{4(1 - \phi)}
\int_{-Nh/2}^{Nh/2} e^{-3x/2}\,\E[L^x_T]\dd x.
\end{equation}
This is obtained from the change of variables $b = e^x$ applied to
\eqref{eq:EGM_limit}, using $A^b_T = L^{\log b}_T/b$ and
$b^{-3/2}\dd b = e^{-3x/2}\dd x$.

Our goal is to compute the ratio $\Psi(\gamma)/F_0^Y(\gamma)$ as an
expansion in powers of $\gamma$ (equivalently, $h$).

%% ===================================================================
\section{The General First-Order Correction}\label{sec:general_correction}

We now state and prove our main structural result: the first-order
correction coefficient for general constant drift.

%% -------------------------------------------------------------------
\subsection{Prefactor expansion}

\begin{lemma}[General prefactor expansion]\label{lem:prefactor_general}
For the prefactor $\psi(h)$ defined in \eqref{eq:psi_def}, with
$\alpha = -2\mu/\sigma^2$:
\begin{equation}\label{eq:psi_expand}
\psi(h) = 1 + h\,\kappa_1^{\psi} + O(h^2),
\end{equation}
where
\begin{equation}\label{eq:kappa1_psi}
\kappa_1^{\psi} = \frac{3}{4} + \frac{\mu}{\sigma^2}
= \frac{3}{4} - \frac{\alpha}{2}.
\end{equation}
\end{lemma}
\begin{proof}
We expand each factor in $c(h)$ separately.

\textbf{Step 1.} The ratio $\gamma/h$:
Since $\gamma = e^h - 1 = h + h^2/2 + h^3/6 + O(h^4)$,
\begin{equation}\label{eq:gamma_over_h}
\frac{\gamma}{h} = 1 + \frac{h}{2} + \frac{h^2}{6} + O(h^3).
\end{equation}

\textbf{Step 2.} The square-root price increment:
\begin{equation}\label{eq:sqrt_expand}
e^{h/2} - 1 = \frac{h}{2} + \frac{h^2}{8} + \frac{h^3}{48}
+ O(h^4).
\end{equation}

\textbf{Step 3.} The crossing probability factor:
\begin{equation}\label{eq:crossing_expand}
\frac{1}{1 + e^{\alpha h}}
= \frac{1}{2 + \alpha h + \alpha^2 h^2/2 + O(h^3)}
= \frac{1}{2}\Bigl(1 - \frac{\alpha h}{2}
+ \frac{\alpha^2 h^2}{4} - \frac{\alpha^2 h^2}{4} + O(h^2)\Bigr)
= \frac{1}{2} - \frac{\alpha h}{4} + O(h^2).
\end{equation}
More precisely:
\begin{equation}\label{eq:crossing_expand_v2}
\frac{1}{1 + e^{\alpha h}}
= \frac{1}{2} - \frac{\alpha h}{4}
+ O(h^3).
\end{equation}
To verify: $1/(2 + u) = (1/2)(1 - u/2 + u^2/4 - \cdots)$ with
$u = \alpha h + \alpha^2 h^2/2 + \cdots$ gives the first-order term
$-\alpha h/4$ as stated.

\textbf{Step 4.} Combining:
\begin{align}
c(h) &= \frac{\gamma}{h} \cdot (e^{h/2} - 1)
\cdot \frac{1}{1 + e^{\alpha h}} \notag \\
&= \Bigl(1 + \frac{h}{2} + O(h^2)\Bigr)
\cdot \Bigl(\frac{h}{2} + \frac{h^2}{8} + O(h^3)\Bigr)
\cdot \Bigl(\frac{1}{2} - \frac{\alpha h}{4} + O(h^2)\Bigr).
\label{eq:c_product}
\end{align}
Collecting terms through order $h^2$:
\begin{align}
c(h) &= \frac{h}{4} \cdot
\Bigl(1 + \frac{h}{2}\Bigr)
\Bigl(1 + \frac{h}{4}\Bigr)
\Bigl(1 - \frac{\alpha h}{2}\Bigr) + O(h^3)
\notag \\
&= \frac{h}{4}\Bigl(1 + \frac{h}{2} + \frac{h}{4}
- \frac{\alpha h}{2} + O(h^2)\Bigr) \notag \\
&= \frac{h}{4}\Bigl(1 + h\bigl(\tfrac{3}{4}
- \tfrac{\alpha}{2}\bigr) + O(h^2)\Bigr).
\label{eq:c_expanded}
\end{align}
Therefore $\psi(h) = 4c(h)/h = (h/4)(4/h)\bigl(1 + h(3/4 - \alpha/2)
+ O(h^2)\bigr)$, but we need to be more careful. From
\eqref{eq:c_expanded}, $4c(h) = h(1 + h(3/4 - \alpha/2) + O(h^2))$,
so
\[
\psi(h) = \frac{4c(h)}{h}
= 1 + h\Bigl(\frac{3}{4} - \frac{\alpha}{2}\Bigr) + O(h^2).
\]
Substituting $\alpha = -2\mu/\sigma^2$ gives
$\kappa_1^{\psi} = 3/4 + \mu/\sigma^2$.
\end{proof}

\begin{remark}\label{rem:gbm_prefactor}
For risk-neutral GBM ($\alpha = 1$):
$\kappa_1^{\psi} = 3/4 - 1/2 = 1/4$. For ABM ($\alpha = 0$):
$\kappa_1^{\psi} = 3/4$.
\end{remark}

%% -------------------------------------------------------------------
\subsection{The weight function and its derivative at the origin}

The crossing-count formula \eqref{eq:Psi_def} can be written as
$\Psi = (\phi/(1{-}\phi)) \cdot (\psi(h)/4) \cdot \sum h\,f(kh)$,
where we define the \emph{crossing weight function}:
\begin{equation}\label{eq:f_def}
f(x) \coloneqq e^{x/2}\,\E[L^x_T].
\end{equation}
Similarly, the EGM asymptotic \eqref{eq:F0Y} involves the
\emph{asymptotic weight function}:
\begin{equation}\label{eq:g_def}
g(x) \coloneqq e^{-3x/2}\,\E[L^x_T].
\end{equation}

The behavior of $f'(0)$ controls the Euler--Maclaurin correction:

\begin{lemma}[Derivative of the weight function]\label{lem:f_prime}
The crossing weight function $f$ satisfies
\begin{equation}\label{eq:f_prime}
f'(0) = \Bigl(\frac{1}{2} + \frac{\mu}{\sigma^2}\Bigr)\,\E[L^0_T]
= \frac{1}{2}(1 - \alpha)\,\E[L^0_T].
\end{equation}
In particular, $f'(0) = 0$ if and only if
$\mu = -\sigma^2/2$ (equivalently, $\alpha = 1$).
\end{lemma}
\begin{proof}
By the product rule:
\begin{equation}\label{eq:f_prime_product}
f'(x) = \frac{1}{2}\,e^{x/2}\,\E[L^x_T]
+ e^{x/2}\,\partial_x \E[L^x_T].
\end{equation}
At $x = 0$: $f'(0) = \tfrac{1}{2}\,\E[L^0_T]
+ \partial_x\E[L^x_T]\big|_{x=0}$.

It remains to compute $\partial_x\E[L^x_T]\big|_{x=0}$.
From \eqref{eq:expected_LT} and \eqref{eq:transition}:
\begin{equation}\label{eq:dLdx}
\partial_x \E[L^x_T] = \sigma^2 \int_0^T
\partial_x p_X(t, 0, x)\dd t.
\end{equation}
The derivative of the Gaussian density is
\begin{equation}\label{eq:dp_dx}
\partial_x p_X(t, 0, x)
= p_X(t, 0, x) \cdot \frac{-(x - \mu t)}{\sigma^2 t}.
\end{equation}
At $x = 0$:
\begin{equation}\label{eq:dp_dx_0}
\partial_x p_X(t, 0, 0) = p_X(t, 0, 0) \cdot \frac{\mu}{\sigma^2},
\end{equation}
since $-(0 - \mu t)/(\sigma^2 t) = \mu/\sigma^2$.
Integrating over $t$:
\begin{equation}\label{eq:dL_dx_0}
\partial_x\E[L^x_T]\big|_{x=0}
= \sigma^2 \cdot \frac{\mu}{\sigma^2} \int_0^T p_X(t,0,0)\dd t
= \frac{\mu}{\sigma^2}\,\E[L^0_T].
\end{equation}
Combining:
\[
f'(0) = \frac{1}{2}\,\E[L^0_T] + \frac{\mu}{\sigma^2}\,\E[L^0_T]
= \Bigl(\frac{1}{2} + \frac{\mu}{\sigma^2}\Bigr)\,\E[L^0_T].
\]
The equivalence with $(1-\alpha)/2$ follows from
$\alpha = -2\mu/\sigma^2$.
\end{proof}

%% -------------------------------------------------------------------
\subsection{Euler--Maclaurin correction from the Riemann sum}

The crossing-count sum $\sum_{k=-N/2}^{N/2-1} h\,f(kh)$ is a
left-endpoint Riemann sum on $[-Nh/2, Nh/2)$ with $N$ terms of
width~$h$. We now compute its deviation from the integral.

\begin{lemma}[Riemann sum correction]\label{lem:EM}
Let $f$ be $C^2$ on $[-a, a]$ with $a = Nh/2$. Then
\begin{equation}\label{eq:EM_formula}
\sum_{k=-N/2}^{N/2-1} h\,f(kh)
= I_f(h) - Nh^2\,\frac{f'(0)}{2} + O(h^3 N^3 \|f''\|_\infty),
\end{equation}
where $I_f(h) = \int_{-Nh/2}^{Nh/2} f(x)\dd x$ and the second term
comes from the boundary evaluation.
\end{lemma}
\begin{proof}
The standard Euler--Maclaurin formula for a left-endpoint Riemann sum
gives
\begin{equation}\label{eq:EM_standard}
\sum_{k=-N/2}^{N/2-1} h\,f(kh)
= I_f(h) + \frac{h}{2}\bigl[f(-\tfrac{Nh}{2})
- f(\tfrac{Nh}{2})\bigr] + O(h^2 \|f'\|_1).
\end{equation}
The boundary evaluation is:
\begin{align}
f(-a) - f(a) &= -2a\,f'(0) + O(a^3\|f'''\|_\infty)
\notag \\
&= -Nh\,f'(0) + O(N^3 h^3).
\label{eq:boundary_eval}
\end{align}
Substituting into \eqref{eq:EM_standard}:
\begin{align}
\sum h\,f(kh) &= I_f(h) + \frac{h}{2}\bigl(-Nh\,f'(0)
+ O(N^3 h^3)\bigr) + O(h^2) \notag \\
&= I_f(h) - \frac{Nh^2}{2}\,f'(0) + O(h^3).
\label{eq:EM_result}
\end{align}
Dividing by $I_f(h) = Nh\,f(0) + O(N^3 h^3)$:
\begin{equation}\label{eq:sum_over_int}
\frac{\sum h\,f(kh)}{I_f(h)}
= 1 - \frac{h\,f'(0)}{2\,f(0)} + O(h^2).
\end{equation}
\end{proof}

This gives the Euler--Maclaurin contribution to $\kappa_1$:
\begin{equation}\label{eq:kappa1_EM}
\kappa_1^{EM} = -\frac{f'(0)}{2\,f(0)}
= -\frac{1}{2}\Bigl(\frac{1}{2} + \frac{\mu}{\sigma^2}\Bigr)
= -\frac{1}{4} - \frac{\mu}{2\sigma^2},
\end{equation}
using $f(0) = \E[L^0_T]$ and Lemma~\ref{lem:f_prime}.

%% -------------------------------------------------------------------
\subsection{The main theorem: general $\kappa_1$}

\begin{theorem}[General first-order correction]\label{thm:general_kappa1}
Let the log-price follow $dX_t = \mu\dd t + \sigma\dd W_t$ with
constant $\mu \in \R$ and $\sigma > 0$. For an $N$-tick LP range
centered at $X_0 = 0$, the expected renormalized fee satisfies
\begin{equation}\label{eq:general_kappa1}
\frac{\Psi(\gamma)}{F_0^Y(\gamma)}
= 1 + \kappa_1(\mu, \sigma)\,\gamma + O(\gamma^2),
\end{equation}
where
\begin{equation}\label{eq:kappa1_formula}
\kappa_1(\mu, \sigma)
= \frac{1}{2} + \frac{\mu}{2\sigma^2}.
\end{equation}
\end{theorem}
\begin{proof}
The ratio $\Psi/F_0^Y$ decomposes as
\begin{equation}\label{eq:decompose}
\frac{\Psi}{F_0^Y}
= \psi(h) \cdot \frac{\sum h\,f(kh)}{I_g(h)},
\end{equation}
where $I_g(h) = \int_{-Nh/2}^{Nh/2} g(x)\dd x$.

\textbf{Step 1: Prefactor.} By
Lemma~\ref{lem:prefactor_general},
$\psi(h) = 1 + h\,\kappa_1^{\psi} + O(h^2)$ with
$\kappa_1^{\psi} = 3/4 + \mu/\sigma^2$.

\textbf{Step 2: Sum-to-integral ratio.} By
Lemma~\ref{lem:EM} and equation \eqref{eq:sum_over_int}:
\[
\frac{\sum h\,f(kh)}{I_f(h)}
= 1 + h\,\kappa_1^{EM} + O(h^2)
\]
with $\kappa_1^{EM} = -1/4 - \mu/(2\sigma^2)$.

\textbf{Step 3: Weight-function ratio.}
We need $I_f(h)/I_g(h)$ to first order. Since $f(x) - g(x) =
(e^{x/2} - e^{-3x/2})\,\E[L^x_T] = (2x + O(x^2))\,\E[L^x_T]$,
the difference integral is
\begin{align}
I_f(h) - I_g(h) &= \int_{-a}^{a}
(e^{x/2} - e^{-3x/2})\,\E[L^x_T]\dd x \notag \\
&= \int_{-a}^a 2x\,\E[L^x_T]\dd x + O(a^3) \notag \\
&= 2\,\E'[L^0_T] \cdot \frac{2a^3}{3} + O(a^5).
\label{eq:If_minus_Ig}
\end{align}
The last line uses $\int_{-a}^a x\,\E[L^x_T]\dd x
= \E'[L^0_T]\cdot 2a^3/3 + O(a^5)$ (integrating by parts or
expanding $\E[L^x_T] = \E[L^0_T] + x\,\E'[L^0_T] + O(x^2)$ and
using the fact that the $\E[L^0_T]$ term integrates to zero by
antisymmetry). This contributes at order $a^3/a = a^2 = O(h^2)$ to
the ratio, which is $O(h^2) = O(\gamma^2)$:
\begin{equation}\label{eq:weight_ratio_first}
\frac{I_f(h)}{I_g(h)} = 1 + O(h^2).
\end{equation}

\textbf{Step 4: Combining.}
\begin{align*}
\frac{\Psi}{F_0^Y}
&= \psi(h) \cdot
\frac{\sum h\,f(kh)}{I_f(h)} \cdot \frac{I_f(h)}{I_g(h)} \\
&= \bigl(1 + h\,\kappa_1^{\psi} + O(h^2)\bigr)
\bigl(1 + h\,\kappa_1^{EM} + O(h^2)\bigr)
\bigl(1 + O(h^2)\bigr) \\
&= 1 + h\bigl(\kappa_1^{\psi} + \kappa_1^{EM}\bigr) + O(h^2).
\end{align*}
Since $h = \gamma + O(\gamma^2)$, replacing $h$ by $\gamma$ at
first order:
\[
\frac{\Psi}{F_0^Y}
= 1 + \gamma\bigl(\kappa_1^{\psi} + \kappa_1^{EM}\bigr) + O(\gamma^2).
\]
Computing:
\[
\kappa_1 = \kappa_1^{\psi} + \kappa_1^{EM}
= \Bigl(\frac{3}{4} + \frac{\mu}{\sigma^2}\Bigr)
+ \Bigl(-\frac{1}{4} - \frac{\mu}{2\sigma^2}\Bigr)
= \frac{1}{2} + \frac{\mu}{2\sigma^2}.
\qedhere
\]
\end{proof}

%% -------------------------------------------------------------------
\subsection{Special cases}

\begin{corollary}[Risk-neutral GBM]\label{cor:gbm}
Under risk-neutral GBM ($\mu = -\sigma^2/2$):
\begin{equation}\label{eq:kappa1_gbm}
\kappa_1 = \frac{1}{2} + \frac{-\sigma^2/2}{2\sigma^2}
= \frac{1}{2} - \frac{1}{4} = \frac{1}{4}.
\end{equation}
\end{corollary}

\begin{corollary}[Physical measure]\label{cor:physical}
Under the physical measure with expected log-return drift
$\mu_P = r - \sigma^2/2$:
\begin{equation}\label{eq:kappa1_physical}
\kappa_1^P = \frac{1}{2} + \frac{r - \sigma^2/2}{2\sigma^2}
= \frac{1}{4} + \frac{r}{2\sigma^2}.
\end{equation}
For assets with positive expected return ($r > 0$), the first-order
correction is strictly larger than the risk-neutral value~$1/4$.
\end{corollary}

\begin{corollary}[Arithmetic Brownian motion]\label{cor:abm}
Under ABM ($\mu = 0$):
\begin{equation}\label{eq:kappa1_abm}
\kappa_1^{ABM} = \frac{1}{2}.
\end{equation}
\end{corollary}

\begin{example}[Bitcoin under physical measure]\label{ex:btc}
For Bitcoin with annualized expected return $r = 0.5$ (50\%) and
volatility $\sigma = 0.8$: $\kappa_1^P = 1/4 + 0.5/(2 \cdot 0.64)
= 0.25 + 0.39 = 0.64$. The physical-measure correction is $2.56$
times the risk-neutral value. For the $1\%$ fee tier
($\gamma = 0.02$), this means a first-order correction of $1.28\%$
rather than $0.50\%$.
\end{example}

\begin{example}[Ethereum under physical measure]\label{ex:eth}
For Ethereum with $r = 0.3$ and $\sigma = 0.7$:
$\kappa_1^P = 0.25 + 0.3/(2 \cdot 0.49) = 0.25 + 0.31 = 0.56$.
The physical-measure correction is $2.22$ times the risk-neutral
value.
\end{example}

%% ===================================================================
\section{The $f'(0) = 0$ Cancellation: A Structural Analysis}
\label{sec:cancellation}

The ``universality'' of $\kappa_1 = 1/4$ in the original
GBM analysis (where it is independent of $\sigma$, $T$, and $N$) is
traced to the vanishing of $f'(0)$. We now provide a structural
explanation of this cancellation and explore its consequences.

%% -------------------------------------------------------------------
\subsection{The cancellation condition}

\begin{proposition}[Characterization of $f'(0) = 0$]
\label{prop:cancellation}
Among constant-drift diffusions \eqref{eq:general_sde}, the crossing
weight function $f(x) = e^{x/2}\,\E[L^x_T]$ satisfies $f'(0) = 0$
if and only if
\begin{equation}\label{eq:cancel_condition}
\frac{\mu}{\sigma^2} = -\frac{1}{2},
\end{equation}
equivalently, $\alpha = 1$, i.e., the drift is the risk-neutral GBM
drift.
\end{proposition}
\begin{proof}
Immediate from Lemma~\ref{lem:f_prime}: $f'(0) = (1/2 + \mu/\sigma^2)
\,\E[L^0_T]$, and $\E[L^0_T] > 0$ for $T > 0$.
\end{proof}

%% -------------------------------------------------------------------
\subsection{Geometric interpretation}

The condition $\mu/\sigma^2 = -1/2$ has a precise geometric meaning.
The weight $e^{x/2}$ in the fee formula arises from the square-root
price $\pi_k = e^{kh/2}$, which grows with rate $+1/2$ in the
exponent. The local time gradient at the origin is
$\partial_x\E[L^x_T]\big|_0 = (\mu/\sigma^2)\,\E[L^0_T]$, which
under risk-neutral GBM gives a rate $-1/2$. These two rates
\emph{exactly cancel}:
\begin{equation}\label{eq:rate_cancel}
\underbrace{+\frac{1}{2}}_{\text{from } e^{x/2}}
+ \underbrace{\frac{\mu}{\sigma^2}}_{\text{from }
\partial_x\E[L^x]} = 0
\quad \iff \quad \mu = -\frac{\sigma^2}{2}.
\end{equation}
Under any other drift, the two rates do not cancel, the function $f$
is not approximately even near the origin, and the Euler--Maclaurin
boundary term contributes at first order.

%% -------------------------------------------------------------------
\subsection{Consequences for LP strategy}

The non-universality of $\kappa_1$ has direct implications for LP
strategy optimization:

\begin{enumerate}[label=(\arabic*)]
\item \textbf{Risk-neutral vs.\ physical measure.} The EGM formulas
are typically derived under the risk-neutral measure (as in our
setup). But an LP's \emph{actual} fee income is determined by the
physical-measure dynamics. If the underlying asset has positive
expected return $r > 0$, the physical-measure correction is larger
than the risk-neutral one. An LP using the risk-neutral correction
to estimate fees will systematically underestimate the finite-tick
effect.

\item \textbf{Directional LPs.} An LP who concentrates liquidity
above the current price (bullish position) effectively faces a
different drift within their range. The correction $\kappa_1$ then
depends on the conditional drift within the range, which may differ
from the unconditional drift.

\item \textbf{Stablecoin pools.} For stablecoin pairs (e.g.,
USDC/USDT), the drift is approximately zero ($\mu \approx 0$), giving
$\kappa_1 \approx 1/2$. This is \emph{twice} the risk-neutral GBM
value, a substantial difference for the narrow-spread pools where
precision matters most.
\end{enumerate}

%% -------------------------------------------------------------------
\subsection{The decomposition table}

Table~\ref{tab:decomposition} summarizes the decomposition of
$\kappa_1$ into its two sources across different models.

\begin{table}[h]
\centering
\caption{Decomposition of $\kappa_1$ into prefactor and
Euler--Maclaurin contributions.}
\label{tab:decomposition}
\begin{tabular}{lcccccc}
\toprule
Model & $\mu$ & $\alpha$ & $\kappa_1^{\psi}$
& $\kappa_1^{EM}$ & $\kappa_1$ & $f'(0)/\E[L^0_T]$ \\
\midrule
GBM ($\PP^*$) & $-\sigma^2/2$ & $1$ & $1/4$ & $0$ & $1/4$ & $0$ \\
ABM & $0$ & $0$ & $3/4$ & $-1/4$ & $1/2$ & $1/2$ \\
Physical ($r > 0$) & $r - \sigma^2/2$ & $1 - 2r/\sigma^2$
& $1/4 + r/\sigma^2$ & $-r/(2\sigma^2)$ & $1/4 + r/(2\sigma^2)$
& $r/(2\sigma^2)$ \\
General & $\mu$ & $-2\mu/\sigma^2$
& $3/4 + \mu/\sigma^2$ & $-1/4 - \mu/(2\sigma^2)$
& $1/2 + \mu/(2\sigma^2)$ & $1/2 + \mu/\sigma^2$ \\
\bottomrule
\end{tabular}
\end{table}

%% ===================================================================
\section{The Full Two-Term Expansion under GBM}\label{sec:gbm}

We now specialize to risk-neutral GBM ($\mu = -\sigma^2/2$,
$\alpha = 1$) to derive the complete two-term expansion including the
second-order, range-dependent correction.

%% -------------------------------------------------------------------
\subsection{Exact crossing-count formula}

\begin{proposition}[Exact formula under GBM]\label{prop:exact}
Under risk-neutral GBM, the expected renormalized $Y$-fee is
\begin{equation}\label{eq:exact}
\Psi(\gamma) = \frac{\phi}{1 - \phi}
\sum_{k=-N/2}^{N/2-1} c(h) \cdot e^{kh/2} \cdot \E^*[L^{kh}_T],
\end{equation}
where $c(h) = \gamma(e^{h/2} - 1)/(h(1 + e^h))$.
\end{proposition}
\begin{proof}
This is equation \eqref{eq:Psi_def} with $\alpha = 1$.
At each visit to grid level $kh$, the walk departs upward with
probability $p_\uparrow = 1/(1 + e^h)$
(Lemma~\ref{lem:p_up} with $\alpha = 1$).
The expected number of visits satisfies
$\E[\text{visits to }kh] = \E[L^{kh}_T]/h + O(1)$ as $h \to 0$
(see \citealt{GobetLandon2014,GobetStazhynski2020}), and the
$O(1)$ remainder contributes $O(\gamma^2)$ to the renormalized fee
after multiplication by the $O(h)$ fee per crossing.
\end{proof}

%% -------------------------------------------------------------------
\subsection{Prefactor under GBM}

\begin{lemma}[GBM prefactor]\label{lem:prefactor_gbm}
Under GBM ($\alpha = 1$):
$\psi(h) = 1 + h/4 + O(h^2)$.
\end{lemma}
\begin{proof}
This is Lemma~\ref{lem:prefactor_general} with $\alpha = 1$:
$\kappa_1^{\psi} = 3/4 - 1/2 = 1/4$.
\end{proof}

%% -------------------------------------------------------------------
\subsection{Euler--Maclaurin analysis under GBM}

Under GBM, $f'(0) = 0$ (Lemma~\ref{lem:f_prime} with $\alpha = 1$).
Therefore the leading Euler--Maclaurin correction vanishes, and we
must go to second order.

\begin{lemma}[GBM Euler--Maclaurin]\label{lem:EM_gbm}
Under GBM, with $f'(0) = 0$:
\begin{equation}\label{eq:EM_gbm}
\sum_{k=-N/2}^{N/2-1} h\,f(kh)
= I_f(h)\bigl(1 + O(h^2)\bigr).
\end{equation}
The boundary correction is $O(h^4)$ because
$f(-a) - f(a) = -2f'(0)\,a + O(a^3) = O(a^3) = O(N^3 h^3)$,
and the $h/2$ prefactor gives $O(N^3 h^4)$.
\end{lemma}

%% -------------------------------------------------------------------
\subsection{Weight-function discrepancy: the second-order term}

\begin{lemma}[Weight-function discrepancy]\label{lem:weight_discrep}
Under GBM ($\alpha = 1$), with $a = Nh/2$:
\begin{equation}\label{eq:weight_ratio}
\frac{I_f(h)}{I_g(h)} = 1 - \frac{N^2 h^2}{6} + O(h^4).
\end{equation}
\end{lemma}
\begin{proof}
\textbf{Step 1: Second derivatives at the origin.}

We need $f''(0)$ and $g''(0)$. Using the notation
$\lambda_0 \coloneqq \E[L^0_T]$ and
$\lambda_2 \coloneqq \partial_x^2\E[L^x_T]\big|_{x=0}$:
\begin{align}
f''(0) &= \frac{1}{4}\,\lambda_0 + \partial_x\E[L^x_T]\big|_0
+ \lambda_2 \notag \\
&= \frac{1}{4}\,\lambda_0 - \frac{1}{2}\,\lambda_0
+ \lambda_2
= -\frac{1}{4}\,\lambda_0 + \lambda_2,
\label{eq:f_second}
\end{align}
where we used $\partial_x\E[L^x_T]\big|_0 = -\lambda_0/2$ under GBM.
Similarly:
\begin{align}
g''(0) &= \frac{9}{4}\,\lambda_0
- 3 \cdot \bigl(-\frac{1}{2}\bigr)\,\lambda_0 + \lambda_2
= \frac{15}{4}\,\lambda_0 + \lambda_2.
\label{eq:g_second}
\end{align}
Here we used $g(x) = e^{-3x/2}\,\E[L^x_T]$, so
$g'(0) = -3\lambda_0/2 + (-\lambda_0/2) = -2\lambda_0$ and
$g''(0) = 9\lambda_0/4 + 2 \cdot (-3/2)(-1/2)\lambda_0 + \lambda_2
= 9\lambda_0/4 + 3\lambda_0/2 + \lambda_2 = 15\lambda_0/4 + \lambda_2$.

\textbf{Step 2: Integral difference.}

Since $f(0) = g(0) = \lambda_0$ and $f'(0) = 0$,
$g'(0) = -2\lambda_0$:
\begin{align}
I_f(h) - I_g(h) &= \int_{-a}^{a}\bigl[f(x) - g(x)\bigr]\dd x
\notag \\
&= \bigl[f'(0) - g'(0)\bigr] \cdot
\underbrace{\int_{-a}^a x\dd x}_{= 0}
+ \frac{f''(0) - g''(0)}{2} \cdot
\int_{-a}^a x^2\dd x + O(a^5).
\label{eq:integral_diff}
\end{align}
Using $\int_{-a}^a x^2\dd x = 2a^3/3$ and
$f''(0) - g''(0) = (-\lambda_0/4 + \lambda_2) -
(15\lambda_0/4 + \lambda_2) = -4\lambda_0$:
\begin{equation}\label{eq:integral_diff_value}
I_f(h) - I_g(h) = \frac{-4\lambda_0}{2} \cdot
\frac{2a^3}{3} + O(a^5) = -\frac{4\lambda_0\,a^3}{3} + O(a^5).
\end{equation}

\textbf{Step 3: Ratio.}

Since $I_g(h) = 2a\,\lambda_0 + O(a^3)$ (the leading term of the
integral of $g$):
\[
\frac{I_f(h)}{I_g(h)} = 1 + \frac{I_f - I_g}{I_g}
= 1 - \frac{4a^3/3}{2a} + O(a^4)
= 1 - \frac{2a^2}{3} + O(a^4).
\]
Substituting $a = Nh/2$:
\[
\frac{I_f(h)}{I_g(h)} = 1 - \frac{N^2 h^2}{6} + O(h^4).
\qedhere
\]
\end{proof}

%% -------------------------------------------------------------------
\subsection{Main theorem under GBM}

\begin{theorem}[Finite-tick correction under GBM]
\label{thm:main}
Under risk-neutral GBM with constant $\sigma > 0$, $T > 0$,
and $N$ even:
\begin{equation}\label{eq:main_thm}
\frac{\Psi(\gamma)}{F_0^Y(\gamma)}
= 1 + \frac{1}{4}\,\gamma - \frac{N^2}{6}\,\gamma^2
+ O(\gamma^3).
\end{equation}
In particular,
$\kappa_1 = \lim_{\gamma \to 0}[\Psi/F_0^Y - 1]/\gamma = 1/4$,
independent of $\sigma$, $T$, and~$N$.
\end{theorem}
\begin{proof}
From Lemma~\ref{lem:prefactor_gbm}, Lemma~\ref{lem:EM_gbm}, and
Lemma~\ref{lem:weight_discrep}:
\begin{align*}
\frac{\Psi}{F_0^Y}
&= \psi(h) \cdot \frac{\sum h\,f(kh)}{I_g(h)} \\
&= \psi(h) \cdot
\frac{\sum h\,f(kh)}{I_f(h)} \cdot \frac{I_f(h)}{I_g(h)} \\
&= \Bigl(1 + \frac{h}{4} + O(h^2)\Bigr)
\Bigl(1 + O(h^2)\Bigr)
\Bigl(1 - \frac{N^2 h^2}{6} + O(h^4)\Bigr) \\
&= 1 + \frac{h}{4} - \frac{N^2 h^2}{6} + O(h^3).
\end{align*}
Since $h = \gamma - \gamma^2/2 + O(\gamma^3)$:
$h/4 = \gamma/4 + O(\gamma^2)$ and
$N^2 h^2/6 = N^2\gamma^2/6 + O(\gamma^3)$.
\end{proof}

\begin{corollary}[Sign crossover]\label{cor:crossover}
The relative correction $\Psi/F_0^Y - 1$ changes sign at
\begin{equation}\label{eq:gamma_star}
\gamma^* = \frac{3}{2N^2}.
\end{equation}
For $\gamma < \gamma^*$, the asymptotic formula underestimates
fees; for $\gamma > \gamma^*$, it overestimates.
\end{corollary}
\begin{proof}
Setting $\gamma/4 - N^2\gamma^2/6 = 0$ gives
$\gamma(1/4 - N^2\gamma/6) = 0$, so
$\gamma^* = 3/(2N^2)$.
\end{proof}

\begin{remark}[Universality of $\kappa_1 = 1/4$: scope]
\label{rem:universality}
Under risk-neutral GBM, $\kappa_1 = 1/4$ is independent of $\sigma$,
$T$, and~$N$. This ``universality'' stems from the algebraic
structure of the prefactor $\psi(h)$, combined with the cancellation
$f'(0) = 0$. However, as shown in
Theorem~\ref{thm:general_kappa1}, it is \emph{not} universal across
drift parameters. The value $1/4$ is a special property of the
risk-neutral GBM drift.
\end{remark}

\begin{remark}[Nature of the second-order term]
The coefficient $-N^2/6$ of the second-order correction encodes a
geometric mismatch between the spatial weight
$e^{x/2}$ (from the fee formula) and the temporal weight
$e^{-3x/2}$ (from the occupation-time representation).
Since $e^{x/2} - e^{-3x/2} = 2x + O(x^2)$, the discrepancy grows
quadratically with the range width $Nh$. The mismatch is always
negative: the crossing-count sum assigns larger weight to ticks above
the center (via $e^{x/2}$), while the local time integral assigns
larger weight below the center (via $e^{-3x/2}$). Since $\E[L^x_T]$
is approximately symmetric near $x = 0$ (the starting point), these
opposing biases produce a deficit in the sum relative to the
integral.
\end{remark}

%% ===================================================================
\section{Misaligned Range Boundaries}\label{sec:misaligned}

In practice, LPs choose their range boundaries $[p_\ell, p_u]$ in
terms of prices, which need not correspond exactly to tick grid
levels. When the boundary falls between two ticks, the crossing
statistics at the boundary tick change. We now quantify this effect.

%% -------------------------------------------------------------------
\subsection{Setup and definitions}

\begin{definition}[Fractional offset]\label{def:offset}
Let $a_\ell$ denote the lower boundary of the LP range in log-price
space. Suppose $a_\ell$ falls between grid levels $mh$ and $(m+1)h$:
\begin{equation}\label{eq:offset}
a_\ell = mh + \delta h, \qquad \delta \in [0, 1).
\end{equation}
The first tick level inside the range is $x_b = (m+1)h$, at distance
\begin{equation}\label{eq:epsilon}
\varepsilon = x_b - a_\ell = (1 - \delta)h
\end{equation}
above $a_\ell$. The case $\delta = 0$ corresponds to alignment
($a_\ell$ exactly at a tick level), while $\delta \to 1$ corresponds
to the range boundary just below a tick level.
\end{definition}

\begin{figure}[h]
\centering
\begin{minipage}{0.9\textwidth}
\small
\begin{verbatim}
        a_l           x_b          x_b+h        x_b+2h
  -------*-------------|-------------|-------------|---
         |<-- eps ---->|
         |<--- (1-d)h -->|<---------- h ---------->|

  Interior ticks: symmetric strip [kh-h, kh+h], width 2h
  Boundary tick:  asymmetric strip [a_l, x_b+h], width (2-d)h
\end{verbatim}
\end{minipage}
\caption{Misaligned range boundary. The boundary tick $x_b$ sits in
an asymmetric strip of width $(2-\delta)h$ instead of the standard
$2h$ strip.}
\label{fig:misaligned}
\end{figure}

%% -------------------------------------------------------------------
\subsection{Boundary tick crossing probability}

\begin{proposition}[Boundary tick efficiency]\label{prop:boundary_eff}
Under the diffusion \eqref{eq:general_sde} with scale function $S$,
the upcrossing probability at the boundary tick $x_b$ (reaching
$x_b + h$ before $a_\ell$) is
\begin{equation}\label{eq:p_boundary}
p_b(\delta, h) = \frac{S(\varepsilon)}{S(h + \varepsilon)}
= \frac{S((1-\delta)h)}{S((2-\delta)h)}.
\end{equation}
Under GBM ($\alpha = 1$):
\begin{equation}\label{eq:p_boundary_gbm}
p_b(\delta, h) = \frac{e^{(1-\delta)h} - 1}{e^{(2-\delta)h} - 1}.
\end{equation}
The boundary tick efficiency ratio, defined as $p_b/p_\uparrow$, satisfies
\begin{equation}\label{eq:R_delta}
R(\delta) \coloneqq \lim_{h \to 0}\frac{p_b(\delta, h)}{p_\uparrow(h)}
= \frac{2(1-\delta)}{2-\delta}.
\end{equation}
\end{proposition}
\begin{proof}
The crossing probability follows from the standard exit problem:
starting at $x_b$, the process must reach $x_b + h$ before $a_\ell$.
The scale function gives
\[
p_b = \frac{S(x_b) - S(a_\ell)}{S(x_b + h) - S(a_\ell)}
= \frac{S(\varepsilon)}{S(\varepsilon + h)},
\]
by translation invariance of constant-coefficient diffusions.

For the efficiency ratio, using the GBM scale function
$S(x) = e^x - 1$:
\begin{align*}
\frac{p_b}{p_\uparrow}
&= \frac{S(\varepsilon)}{S(\varepsilon + h)}
\cdot \frac{S(2h)}{S(h)}
= \frac{(e^{\varepsilon} - 1)(e^{2h} - 1)}{%
(e^{\varepsilon + h} - 1)(e^h - 1)}.
\end{align*}
As $h \to 0$, $S(x) \to x$ (since $\alpha = 1$ and
$(e^x - 1)/x \to 1$), so:
\[
R(\delta) = \frac{\varepsilon}{\varepsilon + h}
\cdot \frac{2h}{h}
= \frac{(1-\delta)h}{(2-\delta)h} \cdot 2
= \frac{2(1-\delta)}{2-\delta}.
\qedhere
\]
\end{proof}

%% -------------------------------------------------------------------
\subsection{Properties of the efficiency function}

\begin{proposition}[Properties of $R(\delta)$]
\label{prop:R_properties}
The boundary tick efficiency $R(\delta) = 2(1-\delta)/(2-\delta)$
satisfies:
\begin{enumerate}[label=(\alph*)]
\item $R(0) = 1$ \textup{(aligned: full efficiency)}.
\item $R(\delta)$ is strictly decreasing on $[0, 1)$.
\item $R(1/2) = 2/3$ \textup{(half-tick offset: $67\%$ efficiency)}.
\item $R(\delta) \to 0$ as $\delta \to 1^-$
\textup{(boundary at tick: zero efficiency)}.
\item $R'(\delta) = -2/(2-\delta)^2 < 0$ for all
$\delta \in [0,1)$.
\item $R(\delta)$ is a \textup{rational} function of $\delta$.
\end{enumerate}
\end{proposition}
\begin{proof}
All properties follow from direct computation:
(a) $R(0) = 2/2 = 1$. (b) and (e): $R'(\delta) = -2/(2-\delta)^2 < 0$.
(c) $R(1/2) = 2 \cdot (1/2)/(3/2) = 2/3$.
(d) $\lim_{\delta \to 1} 2(1-\delta)/(2-\delta) = 0/1 = 0$.
(f) Evident from the formula.
\end{proof}

%% -------------------------------------------------------------------
\subsection{The misalignment correction to fee revenue}

\begin{theorem}[Misalignment correction]\label{thm:misaligned}
For an LP range with fractional offset $\delta$ at the lower boundary
and $\delta'$ at the upper boundary, each with $N$ interior ticks, the
total fee correction acquires an additional first-order term:
\begin{equation}\label{eq:misaligned_correction}
\frac{\Psi_{\mathrm{mis}}}{\Psi_{\mathrm{align}}}
= 1 - \frac{1}{N}\Bigl[\frac{\delta}{2-\delta}
+ \frac{\delta'}{2-\delta'}\Bigr] + O(N^{-2}),
\end{equation}
where $\Psi_{\mathrm{align}}$ is the aligned fee from
Theorem~\ref{thm:main}. The correction is always non-positive,
with equality if and only if both boundaries are aligned
($\delta = \delta' = 0$).
\end{theorem}
\begin{proof}
Under alignment, each tick contributes equally to leading order
(proportional to $\E[L^{kh}_T]$). The boundary ticks each contribute
$p_\uparrow \cdot \E[L^{x_b}_T]/h$ fee events. Under misalignment,
the boundary tick contribution is reduced by the factor $R(\delta)$:
\[
\text{boundary contribution ratio}
= R(\delta) = 1 - \frac{\delta}{2-\delta}.
\]
Since the sum has $N$ terms and each boundary tick is one of them:
\begin{align*}
\Psi_{\mathrm{mis}} &= \Psi_{\mathrm{align}}
\cdot \frac{(N-2) + R(\delta) + R(\delta')}{N} + O(N^{-2}) \\
&= \Psi_{\mathrm{align}} \Bigl(1 - \frac{1}{N}
\bigl[(1 - R(\delta)) + (1 - R(\delta'))\bigr]\Bigr) + O(N^{-2}).
\end{align*}
Using $1 - R(\delta) = \delta/(2-\delta)$ gives
\eqref{eq:misaligned_correction}.
\end{proof}

%% -------------------------------------------------------------------
\subsection{Comparison with the BGK barrier correction}

The classical BGK correction for a discretely monitored barrier option
involves the constant
\begin{equation}\label{eq:bgk_beta}
\beta = \frac{\zeta(1/2)}{\sqrt{2\pi}} \approx -0.5826,
\end{equation}
where $\zeta(1/2) \approx -1.4604$ is the Riemann zeta function at
$1/2$. This constant arises from the Wiener--Hopf factorization of
the distribution of the random walk maximum, via the Spitzer--Baxter
identity.

\begin{remark}[Why BGK does not apply here]\label{rem:bgk}
The BGK correction arises from the \emph{overshoot}: under discrete
temporal monitoring at times $\{0, \Delta t, 2\Delta t, \ldots\}$,
the random walk can jump past a barrier between monitoring dates.
The overshoot distribution determines the barrier shift.

In our spatial setting, the price process $X_t$ is
\emph{continuous}: it crosses every level $x$ with probability one
(by the intermediate value theorem for continuous paths). There is
\emph{no overshoot} past tick levels. The misalignment correction
instead arises from the \emph{gambler's ruin probability} in an
asymmetric strip $[a_\ell, x_b + h]$ of width $(2-\delta)h$ (vs.\
the standard symmetric strip of width $2h$).

Consequently, our correction is a rational function of $\delta$
(Proposition~\ref{prop:R_properties}(f)), not an irrational constant
involving $\zeta(1/2)$. This distinction---spatial discretization
produces algebraic corrections while temporal discretization produces
transcendental ones---appears to be new in the literature.
\end{remark}

\begin{remark}[Sub-tick strip effect]
In addition to the boundary tick efficiency reduction, misalignment
causes the LP to miss fee revenue from the sub-tick strips
$[a_\ell, x_b)$ (below the first tick) and $(x_t, a_u]$ (above the
last tick), where no tick exists to collect fees. These strips account
for a fraction $\delta/(N + 2 - \delta - \delta')$ of the total range
width, producing a direct $O(1/N)$ correction to the EGM integral
denominator $F_0^Y$.
\end{remark}

%% -------------------------------------------------------------------
\subsection{General scale function}

For a general constant-drift diffusion with $\alpha \ne 1$, the
boundary tick efficiency generalizes:
\begin{equation}\label{eq:R_general}
R_\alpha(\delta) = \lim_{h \to 0}
\frac{S_\alpha((1-\delta)h)/S_\alpha((2-\delta)h)}{%
S_\alpha(h)/S_\alpha(2h)}
= \frac{(1-\delta)(2)}{(2-\delta)(1)}
= \frac{2(1-\delta)}{2-\delta},
\end{equation}
which is identical to the GBM case. This is because
$S_\alpha(x)/x \to 1$ as $x \to 0$ for any $\alpha$, and the
efficiency ratio depends only on the linear behavior of $S$ near the
origin. The $\alpha$-dependence enters only at $O(h^2)$ (the $O(h)$ coefficient vanishes by cancellation):
\begin{equation}\label{eq:R_alpha_correction}
\frac{p_b}{p_\uparrow}
= R(\delta) + O(h^2).
\end{equation}

%% ===================================================================
\section{Optimal Range Width}\label{sec:optimal}

The finite-tick correction enables a refined optimization of the LP's
range width, accounting for the systematic bias in the EGM asymptotic.

%% -------------------------------------------------------------------
\subsection{LP revenue decomposition}

The LP's net revenue over $[0, T]$ decomposes as
\begin{equation}\label{eq:net_revenue}
\mathrm{Net}(N, \gamma) = \mathrm{Fee}(N, \gamma)
- \IL(N, \gamma),
\end{equation}
where $\mathrm{Fee}$ is the swap fee income and $\IL$ is the
impermanent loss.

\textbf{Fee revenue.} Under the corrected formula:
\begin{equation}\label{eq:corrected_fee}
\mathrm{Fee}(N, \gamma) = F_0^Y(N, \gamma)
\cdot \Bigl(1 + \frac{1}{4}\,\gamma
- \frac{N^2}{6}\,\gamma^2 + O(\gamma^3)\Bigr).
\end{equation}
For a narrow range ($Nh \ll \sigma\sqrt{T}$), the EGM leading term
scales linearly in $N$:
\begin{equation}\label{eq:F0_narrow}
F_0^Y(N, \gamma) \approx C_{\mathrm{fee}} \cdot N\gamma,
\qquad C_{\mathrm{fee}} = \frac{\phi\,\sigma\sqrt{2T/\pi}}
{4(1 - \phi)}.
\end{equation}
This uses $\int_{-a}^a \E[L^x_T]\dd x \approx 2a\,\E[L^0_T]$ for
$a$ small, together with
$\E[L^0_T] = \sigma^2\int_0^T p_X(t,0,0)\dd t
= \sigma\sqrt{2T/\pi}$.

\textbf{Impermanent loss.} For a narrow range, the IL scales
quadratically:
\begin{equation}\label{eq:IL}
\IL(N, \gamma) \approx C_{\IL} \cdot (N\gamma)^2,
\qquad C_{\IL} = \frac{\sigma^2 T}{24}.
\end{equation}
This is the leading-order expansion of the standard IL formula for
concentrated liquidity positions (see, e.g., \citealt{ClarkHassani2023}).

%% -------------------------------------------------------------------
\subsection{Uncorrected optimum}

Ignoring the finite-tick correction, the LP maximizes
\begin{equation}\label{eq:net_uncorrected}
\mathrm{Net}_0(N) = C_{\mathrm{fee}} \cdot N\gamma
- C_{\IL} \cdot N^2\gamma^2.
\end{equation}
Setting $\partial_N \mathrm{Net}_0 = 0$:
\begin{equation}\label{eq:N_star_EGM}
N^*_{\mathrm{EGM}} = \frac{C_{\mathrm{fee}}}{2\,C_{\IL}\,\gamma}
= \frac{\phi\,\sigma\sqrt{2T/\pi}}{4(1-\phi)}
\cdot \frac{12}{\sigma^2 T\,\gamma}
= \frac{3\phi}{(1-\phi)\,\sigma\,\gamma}
\cdot \sqrt{\frac{2}{\pi T}}.
\end{equation}

%% -------------------------------------------------------------------
\subsection{Corrected optimum}

With the finite-tick correction, the fee becomes
\begin{equation}\label{eq:fee_corrected_full}
\mathrm{Fee}(N, \gamma) \approx C_{\mathrm{fee}} \cdot N\gamma
\Bigl(1 + \frac{\gamma}{4} - \frac{N^2\gamma^2}{6}\Bigr).
\end{equation}
The net revenue is
\begin{equation}\label{eq:net_corrected}
\mathrm{Net}(N) = C_{\mathrm{fee}}\,N\gamma
+ C_{\mathrm{fee}}\,N\,\frac{\gamma^2}{4}
- C_{\mathrm{fee}}\,\frac{N^3\gamma^3}{6}
- C_{\IL}\,N^2\gamma^2.
\end{equation}

\begin{proposition}[Corrected optimal range width]
\label{prop:optimal_N}
The fee-maximizing number of ticks under the corrected formula
satisfies
\begin{equation}\label{eq:N_corrected}
N^*_{\mathrm{corr}} = N^*_{\mathrm{EGM}}
\Bigl(1 - \frac{(N^*_{\mathrm{EGM}})^2\,\gamma^2}{2}
+ O(\gamma)\Bigr).
\end{equation}
The correction is always negative: the finite-tick effect reduces the
optimal range width.
\end{proposition}
\begin{proof}
Setting $\partial_N\,\mathrm{Net} = 0$:
\begin{equation}\label{eq:FOC}
C_{\mathrm{fee}}\,\gamma\Bigl(1 + \frac{\gamma}{4}\Bigr)
- \frac{C_{\mathrm{fee}}\,N^2\gamma^3}{2}
- 2\,C_{\IL}\,N\gamma^2 = 0.
\end{equation}
Write $N = N^*_{\mathrm{EGM}}(1 + \epsilon)$ with
$|\epsilon| \ll 1$. Substituting and linearizing in $\epsilon$:
\[
C_{\mathrm{fee}}\,\gamma \cdot \frac{\gamma}{4}
- \frac{C_{\mathrm{fee}}\,(N^*_{\mathrm{EGM}})^2\gamma^3}{2}
- 2\,C_{\IL}\,N^*_{\mathrm{EGM}}\,\gamma^2\,\epsilon = 0.
\]
Using $2\,C_{\IL}\,N^*_{\mathrm{EGM}}\,\gamma =
C_{\mathrm{fee}}$ from \eqref{eq:N_star_EGM}, we divide by
$C_{\mathrm{fee}}\,\gamma$ to obtain
\[
\epsilon = \frac{\gamma}{4}
- \frac{(N^*_{\mathrm{EGM}})^2\,\gamma^2}{2}.
\]
Since $N^*_{\mathrm{EGM}} = O(\gamma^{-1})$, the product
$(N^*_{\mathrm{EGM}})^2\gamma^2$ is $O(1)$, so the dominant
correction is $\epsilon \approx
-(N^*_{\mathrm{EGM}})^2\gamma^2/2$.
\end{proof}

%% -------------------------------------------------------------------
\subsection{Practical impact}

Table~\ref{tab:optimal_N} summarizes the corrected optimal range
widths for typical parameters.

\begin{table}[h]
\centering
\caption{Optimal range width: uncorrected vs.\ corrected. Parameters:
$\sigma = 0.8$, $T = 1$ day, $\phi = 0.003$.}
\label{tab:optimal_N}
\begin{tabular}{ccccc}
\toprule
Fee tier & $\gamma$ & $N^*_{\mathrm{EGM}}$
& $N^*_{\mathrm{corr}}$ & Tick shift \\
\midrule
$0.05\%$ & $0.001$ & $172$ & $170$ & $-2$ \\
$0.30\%$ & $0.006$ & $29$ & $28$ & $-1$ \\
$1.00\%$ & $0.020$ & $9$ & $9$ & $0$ \\
\bottomrule
\end{tabular}
\end{table}

The correction is most significant for the $0.05\%$ pool where
ranges are widest. The sign crossover condition $N^2\gamma > 3/2$
from Corollary~\ref{cor:crossover} can be rewritten as: an LP
using the uncorrected formula overestimates fees whenever
$N > \sqrt{3/(2\gamma)}$, which corresponds to
$N > 39$ for the $0.05\%$ tier and $N > 9$ for the $1\%$
tier.

For the $0.05\%$ pool with a \$1M LP position, the 2-tick correction
corresponds to approximately \$200/day in optimized fee revenue. Over
a year, this amounts to \$73{,}000---a non-trivial fraction of the
LP's total income.

%% ===================================================================
\section{Numerical Validation}\label{sec:numerics}

We validate Theorems~\ref{thm:general_kappa1} and~\ref{thm:main},
as well as the misalignment correction
(Theorem~\ref{thm:misaligned}), via semi-analytical computation.
The expected local time $\E[L^x_T]$ is evaluated by numerical
quadrature of the Gaussian transition density
\eqref{eq:transition}, and the sums/integrals in $\Psi$ and
$F_0^Y$ are computed directly.

%% -------------------------------------------------------------------
\subsection{General $\kappa_1$ formula}

\begin{table}[h]
\centering
\caption{Verification of $\kappa_1 = 1/2 + \mu/(2\sigma^2)$ across
drift values. Parameters: $\sigma = 0.4$, $T = 1$, $N = 20$.}
\label{tab:general_kappa1}
\begin{tabular}{lccccc}
\toprule
Model & $\mu$ & $\alpha$ & $\kappa_1$ (theory)
& $\kappa_1$ (numerical, $\gamma = 10^{-4}$)
& $\kappa_1$ (numerical, $\gamma = 10^{-3}$) \\
\midrule
GBM ($\PP^*$) & $-0.08$ & $1.00$ & $0.2500$
& $0.236$ & $0.240$ \\
ABM & $0$ & $0$ & $0.5000$
& $0.494$ & $0.497$ \\
$\mu = 0.05$ & $0.05$ & $-0.625$ & $0.6563$
& $0.655$ & $0.656$ \\
$\mu = -0.20$ & $-0.20$ & $2.50$ & $-0.1250$
& $-0.127$ & $-0.126$ \\
$\mu = 0.10$ & $0.10$ & $-1.25$ & $0.8125$
& $0.811$ & $0.812$ \\
\bottomrule
\end{tabular}
\end{table}

The numerical values converge to the theoretical predictions as
$\gamma \to 0$, with residuals consistent with $O(\gamma)$ terms.
The case $\mu = -0.20$ demonstrates that $\kappa_1$ can be negative,
meaning the finite-tick correction reduces rather than increases the
fee estimate for sufficiently negative drift.

%% -------------------------------------------------------------------
\subsection{GBM two-term expansion}

\begin{table}[h]
\centering
\caption{Relative correction under GBM ($\PP^*$).
Parameters: $\sigma = 0.4$, $T = 1$, $\phi = 0.01$, $N = 20$.}
\label{tab:main}
\begin{tabular}{ccccc}
\toprule
$\gamma$ & $\Psi/F_0^Y - 1$ & $\kappa_1(\gamma)$
& $\gamma/4 - N^2\gamma^2/6$ & Residual \\
\midrule
$0.001$ & $+0.016\%$ & $+0.16$ & $+0.018\%$ & $-0.003\%$ \\
$0.005$ & $-0.034\%$ & $-0.07$ & $-0.042\%$ & $+0.008\%$ \\
$0.010$ & $-0.35\%$ & $-0.35$ & $-0.42\%$ & $+0.07\%$ \\
$0.020$ & $-1.63\%$ & $-0.82$ & $-2.17\%$ & $+0.53\%$ \\
\bottomrule
\end{tabular}
\end{table}

The two-term approximation agrees well with the semi-analytical values
for $\gamma \le 0.01$, with residuals under $0.1\%$. At
$\gamma = 0.02$ (the $1\%$ fee tier), the cubic and higher-order
terms contribute approximately $0.5\%$, suggesting that a three-term
expansion would be needed for precise estimates at this coarse tick
spacing.

%% -------------------------------------------------------------------
\subsection{Sensitivity analysis}

\begin{table}[h]
\centering
\caption{Sensitivity to $N$ and $\sigma$ at $\gamma = 0.005$,
$T = 1$ under GBM ($\PP^*$).}
\label{tab:sensitivity}
\begin{tabular}{cccccccc}
\toprule
& \multicolumn{3}{c}{$\sigma = 0.4$ (varying $N$)}
& \phantom{ab}
& \multicolumn{3}{c}{$N = 20$ (varying $\sigma$)} \\
\cmidrule(lr){2-4}\cmidrule(lr){6-8}
$N$ & $\kappa_1^{\mathrm{eff}}$ & Theory & Residual &
& $\sigma$ & $\kappa_1^{\mathrm{eff}}$ & Theory \\
\midrule
$10$ & $+0.17$ & $+0.17$ & $< 0.01$ &
& $0.2$ & $-0.04$ & $-0.08$ \\
$20$ & $-0.07$ & $-0.08$ & $+0.01$ &
& $0.4$ & $-0.07$ & $-0.08$ \\
$40$ & $-0.96$ & $-1.08$ & $+0.12$ &
& $0.6$ & $-0.07$ & $-0.08$ \\
$80$ & $-4.09$ & $-5.08$ & $+0.99$ &
& $0.8$ & $-0.08$ & $-0.08$ \\
\bottomrule
\end{tabular}
\end{table}

\textbf{Range width dependence.} The effective $\kappa_1$ (computed as
$[\Psi/F_0^Y - 1]/\gamma$) deviates from $1/4$ at finite $\gamma$
due to the second-order term $-N^2\gamma^2/6$. The agreement between
theory and numerics is excellent for $N \le 20$ and deteriorates for
$N \ge 40$, where cubic terms become significant.

\textbf{Volatility independence.} The correction is nearly independent
of $\sigma$ (right panel), consistent with the analytic result that
the leading terms $\gamma/4$ and $-N^2\gamma^2/6$ are
$\sigma$-independent under GBM. The mild dependence at small $\sigma$
reflects higher-order terms involving $\sigma$ through the local time
statistics.

%% -------------------------------------------------------------------
\subsection{Sensitivity to time horizon}

\begin{table}[h]
\centering
\caption{Sensitivity to time horizon $T$. Parameters: $\sigma = 0.4$,
$\gamma = 0.005$, $N = 20$, GBM ($\PP^*$).}
\label{tab:T_sensitivity}
\begin{tabular}{cccc}
\toprule
$T$ & $\Psi/F_0^Y - 1$ & $\kappa_1^{\mathrm{eff}}$
& Theory ($\gamma/4 - N^2\gamma^2/6$) \\
\midrule
$0.1$ & $-0.030\%$ & $-0.06$ & $-0.042\%$ \\
$0.5$ & $-0.033\%$ & $-0.07$ & $-0.042\%$ \\
$1.0$ & $-0.034\%$ & $-0.07$ & $-0.042\%$ \\
$5.0$ & $-0.037\%$ & $-0.07$ & $-0.042\%$ \\
\bottomrule
\end{tabular}
\end{table}

The correction is effectively independent of $T$, as predicted by
the theory (all terms in the expansion are $T$-independent). The weak
residual dependence on $T$ enters through higher-order terms in the
Euler--Maclaurin expansion.

%% -------------------------------------------------------------------
\subsection{Misalignment correction}

\begin{table}[h]
\centering
\caption{Boundary tick efficiency ratio $R(\delta)$ under GBM.
Parameters: $\gamma = 0.005$, $N = 20$.}
\label{tab:misalignment}
\begin{tabular}{cccc}
\toprule
$\delta$ & $R(\delta)$ (theory) & $R(\delta)$ (numerical)
& Relative error \\
\midrule
$0.00$ & $1.0000$ & $1.0000$ & $< 10^{-8}$ \\
$0.25$ & $0.8571$ & $0.8571$ & $< 10^{-6}$ \\
$0.50$ & $0.6667$ & $0.6667$ & $< 10^{-6}$ \\
$0.75$ & $0.4000$ & $0.4000$ & $< 10^{-6}$ \\
$0.90$ & $0.1818$ & $0.1818$ & $< 10^{-5}$ \\
$0.95$ & $0.0952$ & $0.0952$ & $< 10^{-5}$ \\
\bottomrule
\end{tabular}
\end{table}

The boundary tick efficiency formula $R(\delta) = 2(1-\delta)/(2-\delta)$
is exact to within numerical precision for all tested offsets.

\begin{table}[h]
\centering
\caption{Impact of misalignment on total fee revenue. Parameters:
$\gamma = 0.005$, GBM ($\PP^*$).}
\label{tab:misalignment_impact}
\begin{tabular}{ccccc}
\toprule
$N$ & $\delta$ & $\Psi_{\mathrm{mis}}/\Psi_{\mathrm{align}}$
& Theory & Residual \\
\midrule
$4$ & $0.50$ & $0.833$ & $0.833$ & $< 0.001$ \\
$10$ & $0.50$ & $0.933$ & $0.933$ & $< 0.001$ \\
$20$ & $0.50$ & $0.967$ & $0.967$ & $< 0.001$ \\
$40$ & $0.50$ & $0.983$ & $0.983$ & $< 0.001$ \\
$20$ & $0.25$ & $0.986$ & $0.986$ & $< 0.001$ \\
$20$ & $0.75$ & $0.940$ & $0.940$ & $< 0.001$ \\
\bottomrule
\end{tabular}
\end{table}

For narrow ranges ($N = 4$), misalignment can reduce fee revenue by
up to $17\%$---far larger than the aligned-range correction. Even for
moderate ranges ($N = 20$), a half-tick offset causes a $3.3\%$
reduction, which is an order of magnitude larger than the aligned
two-term correction at $\gamma = 0.005$.

%% -------------------------------------------------------------------
\subsection{Physical-measure corrections for crypto assets}

\begin{table}[h]
\centering
\caption{First-order correction under the physical measure for
typical crypto parameters. Fee tier: $1\%$ ($\gamma = 0.02$),
$N = 10$.}
\label{tab:crypto}
\begin{tabular}{lccccc}
\toprule
Asset & $r$ & $\sigma$ & $\kappa_1^P$ & Correction
& vs.\ risk-neutral \\
\midrule
BTC & $0.50$ & $0.80$ & $0.641$ & $+1.28\%$
& $\times 2.56$ \\
ETH & $0.30$ & $0.70$ & $0.556$ & $+1.11\%$
& $\times 2.22$ \\
SOL & $0.80$ & $1.10$ & $0.581$ & $+1.16\%$
& $\times 2.32$ \\
Stablecoin & $\approx 0$ & $0.01$ & $\approx 0.50$
& $+1.00\%$ & $\times 2.00$ \\
\bottomrule
\end{tabular}
\end{table}

Under the physical measure, the first-order correction is $2$--$3$
times the risk-neutral value for typical crypto assets. For stablecoin
pools, the correction approaches $1/2$ (the ABM limit), which is
twice the risk-neutral GBM value. These differences are economically
significant: for a \$10M BTC/USDC position in the $1\%$ pool, the
physical-measure correction implies an additional \$12{,}800 per year
in systematic estimation bias relative to the risk-neutral formula.

%% -------------------------------------------------------------------
\subsection{Practical fee-tier summary}

\begin{table}[h]
\centering
\caption{Summary of corrections for standard Uniswap v3 fee tiers
under GBM ($\PP^*$) with typical LP range widths.}
\label{tab:practical}
\begin{tabular}{cccccc}
\toprule
Fee tier & $\gamma$ & $N$ & $\Psi/F_0^Y - 1$
& Direction & \$ impact per \$10M \\
\midrule
$0.01\%$ & $0.0001$ & $100$ & $+0.002\%$ & Under
& $+\$200$ \\
$0.05\%$ & $0.001$ & $20$ & $+0.02\%$ & Under
& $+\$2{,}000$ \\
$0.30\%$ & $0.006$ & $30$ & $-0.35\%$ & Over
& $-\$35{,}000$ \\
$1.00\%$ & $0.020$ & $10$ & $-0.08\%$ & Over
& $-\$8{,}000$ \\
$1.00\%$ & $0.020$ & $20$ & $-1.63\%$ & Over
& $-\$163{,}000$ \\
\bottomrule
\end{tabular}
\end{table}

The correction is most significant for the $1\%$ pool with wide
ranges, and for the $0.30\%$ pool which is the most commonly used
tier. In all cases where $\gamma > \gamma^*$, the EGM asymptotic
\emph{overestimates} fees, meaning LPs using the asymptotic formula
expect more income than they will actually receive.

%% ===================================================================
\section{Discussion}\label{sec:discussion}

%% -------------------------------------------------------------------
\subsection{Economic implications}

\paragraph{Risk-neutral vs.\ physical measure for LP decisions.}
The distinction between $\kappa_1 = 1/4$ (risk-neutral) and
$\kappa_1^P = 1/4 + r/(2\sigma^2)$ (physical) raises a fundamental
question for LP strategy: which measure should LPs use?

Under the risk-neutral measure, the correction formula is
``universal'' ($\sigma$- and $T$-independent at first order) and
analytically convenient. However, the LP's actual P\&L is determined
by the physical-measure dynamics, which includes the expected return.
Since $\kappa_1^P > \kappa_1$ for assets with $r > 0$, an LP using
the risk-neutral correction will \emph{underestimate} the true
finite-tick effect.

For pricing and hedging purposes, the risk-neutral formula is
appropriate. For P\&L prediction and range optimization, the
physical-measure formula should be used. This parallels the standard
distinction in options pricing between risk-neutral and actuarial
valuation.

\paragraph{Implications for range selection.}
The finite-tick correction provides a concrete criterion for when the
EGM asymptotic is reliable: the relative error is below $1\%$ when
$|\gamma/4 - N^2\gamma^2/6| < 0.01$, i.e., when $N < 4/\gamma$ and
$N^2\gamma < 0.06/\gamma$. For the $0.05\%$ tier, this gives
$N < 4{,}000$ (effectively always satisfied); for the $1\%$ tier,
$N < 200$ (also satisfied for typical ranges). The correction becomes
critical when the LP targets a precision better than $1\%$.

\paragraph{Market-making profitability.}
The sign crossover at $\gamma^* = 3/(2N^2)$ has practical
implications: LPs in wide-range positions at coarse tick spacings
(the most capital-efficient strategy) face the largest
overestimation. This means that profitability analyses based on the
EGM formula are systematically optimistic for the most popular LP
strategies.

%% -------------------------------------------------------------------
\subsection{Connection to discrete barrier option corrections}

The analogy between AMM positions and barrier options, noted by
\citet{Paone2025}, extends naturally to the correction structure.
Table~\ref{tab:bgk_comparison} summarizes the key differences.

\begin{table}[h]
\centering
\caption{Comparison of spatial (tick-grid) and temporal
(monitoring-date) discretization corrections.}
\label{tab:bgk_comparison}
\begin{tabular}{lcc}
\toprule
& \textbf{Spatial (this paper)} & \textbf{Temporal (BGK)} \\
\midrule
Grid parameter & $h = \log(1+\gamma)$ & $\Delta t$ \\
Process & Continuous in time & Sampled at $\{k\Delta t\}$ \\
Overshoot & None (a.s.) & Non-trivial \\
Correction type & Algebraic/EM & Wiener--Hopf \\
Key constant & $\kappa_1 = 1/4$ (RN GBM) & $\beta \approx -0.583$ \\
Nature & Rational in offsets & Involves $\zeta(1/2)$ \\
Source & Fee formula + EM & Overshoot distribution \\
\bottomrule
\end{tabular}
\end{table}

The fundamental distinction is topological: spatial discretization
preserves the continuity of sample paths (the process still passes
through every level), while temporal discretization breaks it (the
process can jump past a barrier). The corrections reflect this
difference: our corrections are polynomial/rational in the
discretization parameter, while BGK corrections involve
transcendental constants from the Wiener--Hopf theory of random
walks.

%% -------------------------------------------------------------------
\subsection{Extensions and open problems}

\paragraph{Stochastic volatility.}
For models of the form $dX_t = \mu(V_t)\dd t + \sigma(V_t)\dd W_t$
with stochastic volatility $V_t$, the correction $\kappa_1$ becomes
a functional of the volatility path:
\[
\kappa_1^{\mathrm{SV}}
= \E\Bigl[\frac{1}{2} + \frac{\mu(V_\tau)}{2\sigma(V_\tau)^2}
\Bigl|\, X \text{ crosses tick at time } \tau\Bigr],
\]
where $\tau$ is a typical crossing time. Deriving the precise form
requires the joint distribution of the process and its volatility at
tick-crossing times, which is non-trivial even for the Heston model.

\paragraph{Jump processes.}
For L\'evy-driven log-prices with jumps, the crossing probability
$p_\uparrow$ is no longer given by the scale function formula, and
the Wiener--Hopf factorization becomes relevant even in the spatial
setting (because jumps can overshoot tick levels). The correction is
expected to involve both the BGK-type constant $\beta$ and the
algebraic prefactor correction, with the relative importance depending
on the jump intensity and size distribution.

\paragraph{Multi-pool strategies.}
LPs who provide liquidity across multiple pools (e.g., ETH/USDC at
both $0.05\%$ and $0.30\%$) face a portfolio optimization problem
where the finite-tick corrections at different fee tiers interact.
The correction shifts the optimal capital allocation toward the
finer-tick pools (where the correction is smaller).

\paragraph{Time-varying parameters.}
For non-stationary volatility (which is the norm in crypto markets),
the correction coefficients become time-dependent. A natural extension
would be to derive the correction under a piecewise-constant
volatility model, which approximates the general case via a
time-inhomogeneous diffusion.

\begin{conjecture}[Stochastic volatility correction]
\label{conj:sv}
Under a stochastic volatility model with long-run mean
$\bar{\sigma}^2$ and vol-of-vol $\nu$, the first-order correction
satisfies
\[
\kappa_1^{\mathrm{SV}} = \frac{1}{4}
+ \frac{r}{2\bar{\sigma}^2} + O(\nu^2).
\]
That is, the stochastic volatility correction equals the
constant-volatility formula evaluated at the long-run mean, plus a
second-order correction in the vol-of-vol.
\end{conjecture}

%% ===================================================================
\section{Conclusion}\label{sec:conclusion}

We have established a comprehensive finite-tick correction framework
for concentrated liquidity fee revenue in Uniswap v3/v4. Our three
main contributions are:

\begin{enumerate}[label=(\arabic*)]
\item The first-order correction coefficient
$\kappa_1 = 1/2 + \mu/(2\sigma^2)$ is \emph{not} universal: the
value $1/4$ is specific to risk-neutral GBM and arises from a
precise cancellation $f'(0) = 0$ between the fee formula's
exponential weight and the local time gradient. Under the physical
measure with positive expected return, the correction is $2$--$3$
times larger.

\item The misaligned-range correction is a rational function
$R(\delta) = 2(1-\delta)/(2-\delta)$ of the fractional offset, not a
transcendental constant involving $\zeta(1/2)$ as in the BGK
temporal correction. The distinction arises because continuous price
paths cross every spatial level with probability one (no overshoot).

\item The finite-tick correction shifts the optimal range width
downward by $1$--$2$ ticks, with the largest effect in the finest
fee tier. The correction formula provides an actionable tool for LP
range optimization.
\end{enumerate}

Together, these results transform the EGM asymptotic from a
theoretical limit theorem to a practical fee estimation tool with
quantified accuracy, and reveal the mathematical structure underlying
the interaction between continuous-time finance and discrete protocol
mechanics.

%% ===================================================================

\appendix

\section{Detailed Taylor Expansion of the Prefactor}
\label{app:taylor}

We provide the complete Taylor expansion of $c(h)$ and $\psi(h)$
through second order for general $\alpha$.

\subsection{Expansion of $\gamma/h$}

Since $\gamma = e^h - 1$:
\begin{align}
\frac{\gamma}{h} &= \frac{e^h - 1}{h}
= 1 + \frac{h}{2} + \frac{h^2}{6} + \frac{h^3}{24} + O(h^4).
\label{eq:gamma_h_full}
\end{align}

\subsection{Expansion of $e^{h/2} - 1$}

\begin{align}
e^{h/2} - 1 &= \frac{h}{2} + \frac{h^2}{8} + \frac{h^3}{48}
+ O(h^4). \label{eq:exp_half_full}
\end{align}

\subsection{Expansion of $1/(1 + e^{\alpha h})$}

Let $u = \alpha h + \alpha^2 h^2/2 + \alpha^3 h^3/6 + O(h^4)$.
Then $e^{\alpha h} = 1 + u$ and
\begin{align}
\frac{1}{1 + e^{\alpha h}} &= \frac{1}{2 + u}
= \frac{1}{2}\,\frac{1}{1 + u/2} \notag \\
&= \frac{1}{2}\Bigl(1 - \frac{u}{2} + \frac{u^2}{4}
- O(u^3)\Bigr) \notag \\
&= \frac{1}{2} - \frac{\alpha h}{4}
- \frac{\alpha^2 h^2}{8} + \frac{\alpha^2 h^2}{8}
+ O(h^3) \notag \\
&= \frac{1}{2} - \frac{\alpha h}{4}
+ O(h^3).
\label{eq:crossing_full}
\end{align}
Here we used $u/2 = \alpha h/2 + \alpha^2 h^2/4 + O(h^3)$ and
$(u/2)^2 = \alpha^2 h^2/4 + O(h^3)$.

\subsection{Full expansion of $c(h)$}

Combining through $O(h^3)$ (note that the crossing probability
contributes an $h^3$ term via \eqref{eq:crossing_full}):
\begin{align}
c(h) &= \Bigl(1 + \frac{h}{2} + \frac{h^2}{6} + \frac{h^3}{24}\Bigr)
\Bigl(\frac{h}{2} + \frac{h^2}{8} + \frac{h^3}{48}\Bigr)
\Bigl(\frac{1}{2} - \frac{\alpha h}{4}
+ \frac{\alpha^3 h^3}{48}\Bigr) + O(h^5).
\label{eq:c_full}
\end{align}
Expanding through $O(h^3)$:
\begin{align}
c(h) &= \frac{h}{4}\Bigl[1 + h\Bigl(\frac{3}{4}
- \frac{\alpha}{2}\Bigr)
+ h^2\Bigl(\frac{1}{3} - \frac{3\alpha}{8}\Bigr)
+ O(h^3)\Bigr].
\label{eq:c_full_simplified}
\end{align}
The $h^2$ coefficient $1/3$ (rather than $7/24$) reflects the
additional $h^3$ contribution from each factor that was absent in the
two-term truncation.

Hence:
\begin{equation}\label{eq:psi_full}
\psi(h) = \frac{4c(h)}{h}
= 1 + h\Bigl(\frac{3}{4} - \frac{\alpha}{2}\Bigr)
+ h^2\Bigl(\frac{1}{3} - \frac{3\alpha}{8}\Bigr) + O(h^3).
\end{equation}
For $\alpha = 1$ (GBM): $\psi(h) = 1 + h/4 + h^2(-1/24) + O(h^3)$.

\section{Scale Function Computations}\label{app:scale}

\subsection{General constant-drift diffusion}

For $dX_t = \mu\dd t + \sigma\dd W_t$ with
$\alpha = -2\mu/\sigma^2$:
\begin{equation}\label{eq:scale_general}
S_\alpha(x) = \begin{cases}
\displaystyle\frac{e^{\alpha x} - 1}{\alpha} & \alpha \ne 0, \\[6pt]
x & \alpha = 0.
\end{cases}
\end{equation}

The crossing probability:
\begin{equation}\label{eq:p_up_general}
p_\uparrow(h) = \frac{S_\alpha(h)}{S_\alpha(2h)}
= \frac{e^{\alpha h} - 1}{e^{2\alpha h} - 1}
= \frac{1}{1 + e^{\alpha h}}.
\end{equation}

For specific models:
\begin{center}
\begin{tabular}{lcccc}
\toprule
Model & $\mu$ & $\alpha$ & $S(x)$ & $p_\uparrow(h)$ \\
\midrule
GBM ($\PP^*$) & $-\sigma^2/2$ & $1$ & $e^x - 1$ & $1/(1+e^h)$ \\
ABM & $0$ & $0$ & $x$ & $1/2$ \\
Phys.\ GBM & $r - \sigma^2/2$ & $1 - 2r/\sigma^2$
& $(e^{\alpha x}-1)/\alpha$ & $1/(1+e^{\alpha h})$ \\
\bottomrule
\end{tabular}
\end{center}

\subsection{Boundary efficiency for general $\alpha$}

The boundary crossing probability with offset $\varepsilon = (1-\delta)h$:
\begin{align}
p_b(\delta, h, \alpha)
&= \frac{S_\alpha((1-\delta)h)}{S_\alpha((2-\delta)h)}
= \frac{e^{\alpha(1-\delta)h} - 1}{e^{\alpha(2-\delta)h} - 1}
= \frac{1}{1 + e^{\alpha(1-\delta)h} \cdot
\frac{e^{\alpha h} - 1}{e^{\alpha(1-\delta)h} - 1}}.
\label{eq:p_b_general}
\end{align}

Expanding as $h \to 0$:
\begin{align}
\frac{p_b}{p_\uparrow}
&= \frac{S_\alpha((1-\delta)h)\,S_\alpha(2h)}{%
S_\alpha((2-\delta)h)\,S_\alpha(h)} \notag \\
&= \frac{(1-\delta)h \cdot 2h}{(2-\delta)h \cdot h}
\cdot \bigl(1 + O(\alpha^2 h^2)\bigr)
= \frac{2(1-\delta)}{2-\delta} + O(h^2).
\label{eq:ratio_general}
\end{align}

The $O(h)$ coefficient vanishes identically, and the first
correction enters at $O(h^2)$:
\begin{equation}\label{eq:R_Oh}
\frac{p_b}{p_\uparrow}
= \frac{2(1-\delta)}{2-\delta}
\Bigl(1 + \frac{\alpha^2 h^2\,\delta(1-\delta)}{%
6(2-\delta)} + O(h^3)\Bigr).
\end{equation}
This shows that the leading-order efficiency ratio is
$\alpha$-independent, with the drift entering only at $O(h^2)$.

\section{Proof of the Misalignment Theorem}\label{app:misaligned}

\begin{proof}[Proof of Theorem~\ref{thm:misaligned}]
We provide a detailed proof of the misalignment correction.

\textbf{Step 1: Fee decomposition.}
The aligned fee sum has $N$ terms:
\begin{equation}\label{eq:aligned_sum}
\Psi_{\mathrm{align}} = \frac{\phi}{1-\phi}
\sum_{k=-N/2}^{N/2-1} c(h)\,e^{kh/2}\,\E[L^{kh}_T].
\end{equation}
Under misalignment with offset $\delta$ at the lower boundary and
$\delta'$ at the upper boundary, the sum modifies the first and last
terms:
\begin{align}
\Psi_{\mathrm{mis}} &= \frac{\phi}{1-\phi}\Bigl[
R(\delta)\,c(h)\,e^{x_b h/2}\,\E[L^{x_b h}_T] \notag \\
&\quad + \sum_{k=\text{interior}} c(h)\,e^{kh/2}\,\E[L^{kh}_T]
\notag \\
&\quad + R(\delta')\,c(h)\,e^{x_t h/2}\,\E[L^{x_t h}_T]
\Bigr].
\label{eq:misaligned_sum}
\end{align}

\textbf{Step 2: Ratio.}
To leading order, all terms in the sum contribute equally
(since $\E[L^{kh}_T] \approx \E[L^0_T]$ for $kh$ in the range):
\begin{align}
\frac{\Psi_{\mathrm{mis}}}{\Psi_{\mathrm{align}}}
&\approx \frac{(N-2) + R(\delta) + R(\delta')}{N}
\notag \\
&= 1 - \frac{(1-R(\delta)) + (1-R(\delta'))}{N}.
\label{eq:ratio_derivation}
\end{align}

\textbf{Step 3: Substitution.}
Using $1 - R(\delta) = 1 - 2(1-\delta)/(2-\delta)
= \delta/(2-\delta)$:
\begin{equation}\label{eq:final_correction}
\frac{\Psi_{\mathrm{mis}}}{\Psi_{\mathrm{align}}}
= 1 - \frac{1}{N}\Bigl[\frac{\delta}{2-\delta}
+ \frac{\delta'}{2-\delta'}\Bigr] + O(N^{-2}).
\end{equation}
The $O(N^{-2})$ remainder comes from the variation of $\E[L^{kh}_T]$
across the range (the boundary ticks have slightly different local
time values than the center).
\end{proof}

\section{Second-Order Expansion Details}\label{app:second_order}

We collect the detailed computations for the second derivatives of
$f$ and $g$ used in Lemma~\ref{lem:weight_discrep}.

\subsection{Second derivative of $\E[L^x_T]$}

From \eqref{eq:dp_dx}:
\begin{align}
\partial_x^2 p_X(t, 0, x)
&= p_X \cdot \Bigl[\frac{(x-\mu t)^2}{\sigma^4 t^2}
- \frac{1}{\sigma^2 t}\Bigr].
\label{eq:d2p}
\end{align}
At $x = 0$:
\begin{equation}\label{eq:d2p_0}
\partial_x^2 p_X(t, 0, 0)
= p_X(t, 0, 0)\Bigl[\frac{\mu^2}{\sigma^4}
- \frac{1}{\sigma^2 t}\Bigr].
\end{equation}
Integrating:
\begin{equation}\label{eq:d2L_0}
\partial_x^2\E[L^x_T]\big|_0
= \frac{\mu^2}{\sigma^4}\,\E[L^0_T]
- \sigma^2\int_0^T \frac{p_X(t,0,0)}{\sigma^2 t}\dd t
= \frac{\mu^2}{\sigma^4}\,\E[L^0_T] - \lambda_{-1},
\end{equation}
where $\lambda_{-1} = \int_0^T p_X(t,0,0)/t\dd t$ is a divergent
integral that cancels in the difference $f''(0) - g''(0)$.

\subsection{Second derivative of $f$}

From $f(x) = e^{x/2}\,\E[L^x_T]$:
\begin{align}
f''(x) &= \frac{1}{4}\,e^{x/2}\,\E[L^x_T]
+ e^{x/2}\,\partial_x\E[L^x_T]
+ e^{x/2}\,\partial_x^2\E[L^x_T].
\label{eq:f_second_full}
\end{align}
At $x = 0$:
\begin{equation}\label{eq:f_second_0}
f''(0) = \frac{1}{4}\,\lambda_0
+ \frac{\mu}{\sigma^2}\,\lambda_0 + \lambda_2,
\end{equation}
where $\lambda_0 = \E[L^0_T]$ and
$\lambda_2 = \partial_x^2\E[L^x_T]\big|_0$.

Under GBM ($\mu = -\sigma^2/2$):
$f''(0) = \lambda_0/4 - \lambda_0/2 + \lambda_2
= -\lambda_0/4 + \lambda_2$.

\subsection{Second derivative of $g$}

From $g(x) = e^{-3x/2}\,\E[L^x_T]$:
\begin{equation}\label{eq:g_second_0}
g''(0) = \frac{9}{4}\,\lambda_0
- \frac{3\mu}{\sigma^2}\,\lambda_0 + \lambda_2.
\end{equation}
Under GBM:
$g''(0) = 9\lambda_0/4 + 3\lambda_0/2 + \lambda_2
= 15\lambda_0/4 + \lambda_2$.

\subsection{Difference}

Under GBM:
\begin{equation}\label{eq:fg_second_diff}
f''(0) - g''(0) = (-\lambda_0/4 + \lambda_2)
- (15\lambda_0/4 + \lambda_2) = -4\lambda_0.
\end{equation}
Note that $\lambda_2$ cancels. This cancellation is essential: the
second-order correction $-N^2 h^2/6$ is determined entirely by
$\lambda_0$, not by the more complex quantity $\lambda_2$.

For general $\alpha$:
\begin{align}
f''(0) - g''(0)
&= \Bigl(\frac{1}{4} + \frac{\mu}{\sigma^2}\Bigr)\lambda_0
- \Bigl(\frac{9}{4} - \frac{3\mu}{\sigma^2}\Bigr)\lambda_0
\notag \\
&= \Bigl(-2 + \frac{4\mu}{\sigma^2}\Bigr)\lambda_0
= -(2 + 2\alpha)\lambda_0.
\label{eq:fg_diff_general}
\end{align}
Under GBM ($\alpha = 1$): $f''(0) - g''(0) = -4\lambda_0$, as
before.

%% ===================================================================
\bibliographystyle{plainnat}
\begin{thebibliography}{99}

\bibitem[Adams et~al.(2021)]{Adams2021}
H.~Adams, N.~Zinsmeister, M.~Salem, R.~Keefer, and D.~Robinson.
\newblock Uniswap v3 core.
\newblock Technical report, Uniswap Labs, 2021.

\bibitem[Broadie et~al.(1997)]{BGK1997}
M.~Broadie, P.~Glasserman, and S.~Kou.
\newblock A continuity correction for discrete barrier options.
\newblock \emph{Math.\ Finance}, 7(4):325--349, 1997.

\bibitem[Broadie et~al.(1999)]{BGK1999}
M.~Broadie, P.~Glasserman, and S.~Kou.
\newblock Connecting discrete and continuous path-dependent options.
\newblock \emph{Finance Stoch.}, 3:55--82, 1999.

\bibitem[Carta and Fabrizio(2024)]{CartaFabrizio2024}
A.~Carta and S.~Fabrizio.
\newblock Fee revenue for concentrated liquidity providers in
  decentralized exchanges.
\newblock \emph{Quant.\ Finance}, 24(6):789--805, 2024.

\bibitem[Clark and Hassani(2023)]{ClarkHassani2023}
A.~Clark and B.~K.~Hassani.
\newblock Liquidity provider returns in geometric mean markets.
\newblock \emph{Digital Finance}, 5:9--45, 2023.

\bibitem[Echenim et~al.(2025)]{EGM2025}
M.~Echenim, E.~Gobet, and A.-C.~Maurice.
\newblock Uniswap v3: impermanent loss modeling and swap fees asymptotic
  analysis.
\newblock \emph{SIAM J.\ Financial Math.}, 16(2):310--349, 2025.

\bibitem[Feng and Linetsky(2008)]{FengLinetsky2008}
L.~Feng and V.~Linetsky.
\newblock Pricing discretely monitored barrier options and defaultable
  bonds in {L}\'evy process models: a fast {H}ilbert transform approach.
\newblock \emph{Math.\ Finance}, 18(3):337--384, 2008.

\bibitem[Fukasawa et~al.(2023)]{FukasawaMajoWorhacz2023}
M.~Fukasawa, B.~Maj\'o, and M.~Worhacz.
\newblock Optimal liquidity provision in concentrated liquidity markets.
\newblock Working paper, 2023.

\bibitem[Gobet and Landon(2014)]{GobetLandon2014}
E.~Gobet and N.~Landon.
\newblock Almost sure optimal hedging strategy.
\newblock \emph{Ann.\ Appl.\ Probab.}, 24(4):1652--1690, 2014.

\bibitem[Gobet and Stazhynski(2020)]{GobetStazhynski2020}
E.~Gobet and U.~Stazhynski.
\newblock Parametric inference for diffusions observed at stopping times.
\newblock \emph{Electron.\ J.\ Statist.}, 14(1):2098--2122, 2020.

\bibitem[H\"{o}rfelt(2003)]{Horfelt2003}
P.~H\"{o}rfelt.
\newblock Extension of the corrected barrier approximation by
  {B}roadie, {G}lasserman, and {K}ou.
\newblock \emph{Finance Stoch.}, 7(2):231--243, 2003.

\bibitem[Howison and Steinberg(2007)]{HowisonSteinberg2007}
S.~Howison and M.~Steinberg.
\newblock A matched asymptotic expansions approach to continuity
  corrections for discretely sampled options. {P}art 1: barrier options.
\newblock \emph{Appl.\ Math.\ Finance}, 14(1):63--89, 2007.

\bibitem[Paone(2025)]{Paone2025}
A.~Paone.
\newblock Designing barrier options using concentrated liquidity of
  constant function market makers.
\newblock SSRN 5233590, 2025.

\end{thebibliography}

\end{document}
